\documentclass[aps, prx, nofootinbib, superscriptaddress]{revtex4-2}

\usepackage[utf8]{inputenc}
\usepackage{graphicx}
\usepackage{amsmath, amssymb, amsfonts, amsthm, bm, nicefrac}
\usepackage{multirow} 
\usepackage{tabularx} 
\usepackage{array}
\usepackage{units}
\usepackage{tensor} 
\usepackage{braket}
\usepackage{hyperref}
\usepackage{bm}
\usepackage{enumitem}
\usepackage{array}
\setlist[itemize,2]{label=\textopenbullet} % itemize의 두 번째 레벨의 기호를 \textbullet로 설정
\usepackage{booktabs} % 테이블 선을 위한 패키지 추가


\usepackage{xcolor}
\usepackage{soul}
\usepackage[export]{adjustbox}
\usepackage{tikz}
\usetikzlibrary{positioning}



% Define global settings and functions before \begin{document}
\definecolor{SkyBlue}{rgb}{0.5, 0.8, 1} 
\definecolor{MyPink}{rgb}{1, 0.682, 0.843}
\definecolor{customblue}{HTML}{1f77b4}
\definecolor{customorange}{HTML}{ff7f0e}
\definecolor{customgreen}{HTML}{2ca02c}
\definecolor{custompurple}{HTML}{C4A3F0}





% Define a custom function for magnetic ordering globally
\newcommand{\magneticOrdering}[5]{
    % Arguments: { #1: grid row }{ #2: grid column }{ #3: angle theta in degrees }{ #4: length of the arrow }{ #5: color of the arrow }

    % Calculate the grid point's actual position
    \pgfmathsetmacro\xa{#2 * \size - #1 * 0.5 * \size}
    \pgfmathsetmacro\ya{#1 * \size * 0.866} % Height of equilateral triangle

    % Calculate arrow end point
    \pgfmathsetmacro\dx{#4 * cos(#3)} % x-component of the arrow
    \pgfmathsetmacro\dy{#4 * sin(#3)} % y-component of the arrow

    % Draw arrow
    \draw[->, line width=1.5pt] (\xa, \ya) -- ++(\dx, \dy);

    % Draw red dot
    \fill[#5] (\xa, \ya) circle(0.2);
}



\setcounter{secnumdepth}{2}
\makeatletter 
\renewcommand\twocolumngrid{
	\def\footnoterule{
		\dimen@\skip\footins\divide\dimen@\thr@@
		\kern-\dimen@\hrule width.5in\kern\dimen@}
	\do@columngrid{mlt}{\tw@}
 }

\makeatother  
\newtheorem{theorem}{Theorem}[section]
\newtheorem{lemma}[theorem]{Lemma}

\newtheorem{definition}{Definition}[section]
\newtheorem{prop}{Proposition}
\newtheorem{axiom}{Axiom}
\newtheorem{claim}{Claim}
\newtheorem{corollary}{Corollary}[theorem]

\hypersetup{
    colorlinks=true,
    linkcolor=blue,
    filecolor=magenta,      
    urlcolor=blue,                                         %% 
    pdftitle={Stability of Chirality for mixed states}, %% PDF의 제목
    pdfpagemode=FullScreen,
    }

\newcommand{\SP}[1]{{\color{cyan}\footnotesize{(SP) #1}}}
\newcommand{\NBCPatom}{$\mathrm{Na_2 Ba Co (PO_4)_2}$}
\newcommand{\NBCP}{{\sf NBCP}}

\newcommand{\kvec}{\mathbf{k}}
\newcommand{\bmdelta}{\bm{\delta}}

\newcommand{\ha}{\hat{a}}
\newcommand{\hc}{\hat{c}}
\newcommand{\hI}{\hat{\mathbb{I}}}
\newcommand{\hP}{\hat{P}}

\newcommand{\rD}{\mathrm{D}}
\newcommand{\rU}{\mathrm{U}}
\newcommand{\ri}{\mathrm{i}}


\newcommand{\bfr}{\mathbf{r}}




\newcommand{\cA}{\mathcal{A}}
\newcommand{\cB}{\mathcal{B}}
\newcommand{\cC}{\mathcal{C}}
\newcommand{\cD}{\mathcal{D}}
\newcommand{\cE}{\mathcal{E}}
\newcommand{\cG}{\mathcal{G}}
\newcommand{\cH}{\mathcal{H}}
\newcommand{\cJ}{\mathcal{J}}
\newcommand{\cN}{\mathcal{N}}
\newcommand{\cM}{\mathcal{M}}
\newcommand{\cP}{\mathcal{P}}
\newcommand{\cR}{\mathcal{R}}
\newcommand{\cQ}{\mathcal{Q}}
\newcommand{\cU}{\mathcal{U}}


\newcommand{\Id}{\mathrm{Id}}

\newcommand{\bth}{\bm{\theta}}


\newcommand{\Tr}{\mathrm{Tr}}
\newcommand{\Pf}{\mathrm{Pf}}


\newcommand{\Theme}[1]{{\vspace{0.5cm}\noindent\textbf{#1}}}

\begin{document}

\title{Note: Linear Spin Wave Theory}
\author{Sung-Min Park}
\begin{abstract}
    In this article, we summarize the linear spin wave theory (LSWT) described in the literature.
\end{abstract}
\maketitle

\tableofcontents

\newpage
\section*{Summary of Notation and Symbols}



We summarize the symbols and notation used in this notes.
\begin{table*}[t]
    \centering
    \caption{Summary of Notation and Symbols in Linear Spin Wave Theory}
    \begin{tabular}{c|l}
    \multicolumn{2}{c}{\textbf{Spin Hamiltonian and Related Quantities}} \\
    \hline
    $\hat{H}$ & Spin Hamiltonian \\
    $J_{ij}^{\alpha\beta}$ & Exchange coupling between sites $i$ and $j$ for spin components $\alpha$ and $\beta$ \\
    $h_j^{\alpha}$ & External magnetic field at site $j$ along direction $\alpha$ \\
    $E_{\rm cl}$ & Classical ground state energy \\
    $\mathbf{S}_j$ & Classical spin vector at site $j$ \\
    $\hat{\mathbf{S}}_j$ & Quantum spin operator at site $j$ \\
    $\delta\hat{\mathbf{S}}_j$ & Quantum fluctuation operator ($\hat{\mathbf{S}}_j = \mathbf{S}_j + \delta\hat{\mathbf{S}}_j$) \\
    $\widetilde{\mathbf{S}}_j$ & Spin operator in the local reference frame \\
    \hline
    \multicolumn{2}{c}{\textbf{Index Conventions and Coordinate Systems}} \\
    \hline
    $i, j$ & Unit cell indices \\
    $\mu, \nu, \sigma, \rho$ & Sublattice indices (typically $a, b, c, \dots$) \\
    $I = (i, \mu), J = (j, \nu)$ & Composite site indices combining unit cell and sublattice \\
    $\mathbf{r}_i$ & Position of unit cell $i$ \\
    $\boldsymbol{\delta}_\mu$ & Position of sublattice $\mu$ within the unit cell \\
    $\mathbf{R}_I = \mathbf{r}_i + \boldsymbol{\delta}_\mu$ & Position of spin at site $I = (i, \mu)$ \\
    $\alpha, \beta$ & Spin component indices ($x, y, z$ or $+, -, 0$) \\
    $\mathbf{R}_j$ & Rotation matrix from laboratory frame to local frame at site $j$ \\
    $\hat{\mathbf{e}}_j^{\alpha}$ & Basis vectors of local frame at site $j$ ($\alpha = +, -, 0$) \\
    \hline
    \multicolumn{2}{c}{\textbf{Boson Operators and Transformations}} \\
    \hline
    $\hat{a}_j, \hat{a}_j^{\dagger}$ & Bosonic annihilation and creation operators \\
    $\hat{n}_j = \hat{a}_j^{\dagger}\hat{a}_j$ & Boson number operator \\
    $\hat{b}_{\mathbf{k}\mu}, \hat{b}_{\mathbf{k}\mu}^{\dagger}$ & Momentum-space bosonic operators for sublattice $\mu$ \\
    $\hat{\beta}_{\mathbf{k}\mu}, \hat{\beta}_{\mathbf{k}\mu}^{\dagger}$ & Diagonalized magnon operators (after Bogoliubov transformation) \\
    $\Psi_{\mathbf{k}} = (\hat{\mathbf{b}}_{\mathbf{k}}, \hat{\mathbf{b}}_{-\mathbf{k}}^{\dagger})^T$ & Nambu spinor in momentum space \\
    $\widetilde{\Psi}_{\mathbf{k}} = (\hat{\boldsymbol{\beta}}_{\mathbf{k}}, \hat{\boldsymbol{\beta}}_{-\mathbf{k}}^{\dagger})^T$ & Diagonalized Nambu spinor \\
    \hline
    \multicolumn{2}{c}{\textbf{Matrices and Transformation Operators}} \\
    \hline
    $\mathsf{H}_{\mathbf{k}}$ & Hamiltonian matrix in BdG form \\
    $\mathsf{A}_{\mathbf{k}}, \mathsf{B}_{\mathbf{k}}$ & Blocks of the Hamiltonian matrix $\mathsf{H}_{\mathbf{k}}$ \\
    $\mathsf{T}_{\mathbf{k}}$ & Paraunitary transformation matrix for Bogoliubov diagonalization \\
    $\mathsf{E}_{\mathbf{k}}$ & Diagonal matrix of magnon energies \\
    $\mathsf{N}_{\mathbf{k}}(t)$ & Time-dependent correlation matrix in magnon basis \\
    $\mathsf{U}^{\alpha}$ & Spin projection operator in laboratory frame \\
    $\mathsf{S}_{\mathbf{k}}$ & Matrix encoding spin magnitudes and sublattice phase factors \\
    $\mathsf{C}^{\alpha\beta}(\mathbf{k}, t)$ & Correlation matrix in laboratory frame \\
    $\mathsf{J}$ & Symplectic metric tensor \\
    \hline
    \multicolumn{2}{c}{\textbf{Physical Quantities}} \\
    \hline
    $E_{\mathbf{k},\mu}$ & Magnon energy for band $\mu$ at momentum $\mathbf{k}$ \\
    $n_{\mathbf{k},\mu} = 1/(e^{\beta E_{\mathbf{k},\mu}}-1)$ & Bose-Einstein distribution function \\
    $E_0$ & Zero-point energy \\
    $\mathcal{C}^{\alpha\beta}_{\mu\nu}(\mathbf{k}, t)$ & Sublattice spin-spin correlation function \\
    $\mathcal{S}^{\alpha\beta}(\mathbf{k}, \omega)$ & Dynamic structure factor \\
    $\mathcal{A}^{\alpha\beta}(\mathbf{k}, \omega)$ & Spectral function \\
    $\Omega_{n,\mathbf{k}}$ & Berry curvature of the $n$-th magnon band \\
    $C_n$ & Chern number of the $n$-th magnon band \\
    \hline
    \end{tabular}
\end{table*}


\newpage
\section{Introduction to Spin Wave Theory (SWT)}

For the spin wave theory, the spin model of our interest is 
\begin{equation}\label{eq: General model Hamiltonian for spin wave theory}
    \hat{H} 
    = \sum_{ij} \sum_{\alpha, \beta} \hat{S}^{\alpha}_{i} J_{ij}^{\alpha\beta}  \hat{S}^{\beta}_{j}  
    -  \sum_{j, \alpha} h_{j}^{\alpha} \hat{S}^{\alpha}_{j},
\end{equation}
where ${\bf S}_{i} = (S^{x}_{i}, S^{y}_{i}, S_{i}^{z})$ is the spin operator at site $i$, with $S^\alpha_i$ representing the $\alpha$-direction component, and $ij$ denotes the \emph{link}, not the lattice indices\footnote{
    If we count $i$ and $j$ as independent variables, we change the exchange Hamiltonian: $\sum_{ij} \sum_{\alpha, \beta} S^{\alpha}_{i} J_{ij}^{\alpha\beta}  S^{\beta}_{j} \rightarrow \frac{1}{2} \sum_{ij} \sum_{\alpha, \beta} S^{\alpha}_{i} J_{ij}^{\alpha\beta} S^{\beta}_{j}$.
}.
The first term is called the exchange Hamiltonian, and the second term is the Zeeman term. 
Note that many Mott insulators can be generally modeled as in Eqn.~\eqref{eq: General model Hamiltonian for spin wave theory}.
\begin{enumerate}[label=(\roman*)]
    \item {\bf Spin-Exchange interactions:} 
    \begin{equation}
        \sum_{ij} \sum_{\alpha, \beta} \hat{S}^{\alpha}_{i} J_{ij}^{\alpha\beta}  \hat{S}^{\beta}_{j}  
        = \sum_{i \neq j} \hat{\bf S}_{i} \cdot {\bf J}_{ij} \cdot \hat{\bf S}_{j} + \sum_{i} \hat{\bf S}_{i} \cdot {\bf A}_{i} \cdot \hat{\bf S}_{i}
    \end{equation}
    where the exchange energy is composed of inter-site interactions (first term) and single-ion anisotropies (second term). The exchange matrix ${\bf J}_{ij}$ can represent various types of interactions including isotropic exchange, Dzyaloshinskii-Moriya interactions, and anisotropic exchanges such as Kitaev interactions. The single-ion anisotropy term, for example $- A \sum_{i} (S^{x}_{i})^{2}$, describes easy-axis anisotropy along $x$.
    
    \item {\bf External magnetic fields and $g$-tensors:}
    \begin{equation}
        \sum_{j} {\bf h}_{j} \cdot \hat{\bf S}_{j} = \mu_B \sum_{j} \sum_{\alpha, \beta} B_{j}^{\alpha} g_{j}^{\alpha\beta} \hat{S}^{\beta}_{j}
    \end{equation}
    where $\mu_B = 0.057883$ meV/T is the Bohr magneton. In magnetic materials, the $g$-tensor $g_{j}^{\alpha\beta}$ mediates the interaction between spins and external magnetic fields and can exhibit significant anisotropy due to spin-orbit coupling and crystal field effects, unlike the isotropic $g$-factor of free electrons ($g_e \approx 2.002$).
\end{enumerate}

\subsection{Spin to Boson Transformations}\label{subsec: spin-to-boson transformation}
A spin wave theory for the model [Eqn.~\eqref{eq: General model Hamiltonian for spin wave theory}] can be formulated by assuming the quantum spin exhibits small fluctuations around its classical spin configuration:
\begin{equation}
    \hat{\mathbf{S}}_{j} = \mathbf{S}_{j} + \delta \hat{\mathbf{S}}_{j},
\end{equation}
where $\mathbf{S}_{j}$ is the classical spin and the operator $\delta \hat{\mathbf{S}}_{j}$ represents the quantum fluctuation. Substituting this into the Hamiltonian $H$, we obtain:
\begin{align*}
    \hat{H} 
    = & \sum_{ij} \big( \mathbf{S}_{i} + \delta \hat{\mathbf{S}}_{i} \big) \cdot \mathbf{J}_{ij} \cdot \big( \mathbf{S}_{j} + \delta \hat{\mathbf{S}}_{j} \big) - \sum_{i} \mathbf{h}_{i} \cdot \big( \mathbf{S}_{i} + \delta \hat{\mathbf{S}}_{i} \big) \\
    = & E_{\rm cl}
    + \Big[ \sum_{ij} \big( \delta \hat{\mathbf{S}}_{i} \cdot \mathbf{J}_{ij} \cdot \mathbf{S}_{j} + \mathbf{S}_{i} \cdot \mathbf{J}_{ij} \cdot \delta \hat{\mathbf{S}}_{j} \big) - \sum_{i} \mathbf{h}_{i} \cdot \delta \hat{\mathbf{S}}_{i} \Big] 
    + \sum_{ij} \delta \hat{\mathbf{S}}_{i} \cdot \mathbf{J}_{ij} \cdot \delta \hat{\mathbf{S}}_{j} \\
    = & E_{\rm cl}
    + \sum_{i} \Big( \sum_{j \in \mathsf{L}(i)} \mathbf{S}_{j} \cdot \mathbf{J}_{ji} - \mathbf{h}_{i} \Big) \cdot \delta \hat{\mathbf{S}}_{i}
    + \sum_{i,j} \delta \hat{\mathbf{S}}_{i} \cdot \mathbf{J}_{ij} \cdot \delta \hat{\mathbf{S}}_{j},
\end{align*}
where $E_{\rm cl} = \sum_{i,j} \mathbf{S}_{i} \cdot \mathbf{J}_{ij} \cdot \mathbf{S}_{j} - \sum_{i} \mathbf{h}_{i} \cdot \mathbf{S}_{i}$ is the classical energy, and the summation runs over the links, not the lattice indices\footnote{
    If we count the indices independently, this becomes $\sum_{i} \Big( 2 \sum_{j} \mathbf{S}_{j} \cdot \mathbf{J}_{ji} - \mathbf{h}_{i} \Big) \cdot \delta \hat{\mathbf{S}}_{i}$.
    },
and $\sum_{i} \sum_{j \in \mathsf{L}(i)} = \sum_{\langle i, j \rangle}$.
In the second line, we use $\delta \hat{\mathbf{S}}_{i} \cdot \mathbf{J}_{ij} \cdot \mathbf{S}_{j} = \mathbf{S}_{j} \cdot \mathbf{J}_{ji} \cdot \delta \hat{\mathbf{S}}_{i}$, which follows from $J_{ij}^{\alpha\beta} = J_{ji}^{\beta\alpha}$.



The fluctuation operator $\delta \hat{\mathbf{S}}_{j}$ can be further decomposed into components perpendicular and parallel to the classical spin orientation:
\begin{equation}
    \delta \hat{\mathbf{S}}_{j} = \delta \hat{\mathbf{S}}_{j}^{\perp} + \delta \hat{\mathbf{S}}_{j}^{\parallel} 
    %= \big( \widetilde{S}_{j}^{+} \hat{\bf e}_{j}^{-} + \widetilde{S}_{j}^{-} \hat{\bf e}_{j}^{+} \big) + \widetilde{S}^{z}_{j} \hat{\bf e}_{j}^{0}.
\end{equation}
Formally, this can be done using the \emph{Holstein-Primakoff} transformation~\cite{PhysRev.58.1098}, which represents the spin operator in the laboratory frame in terms of bosonic operators:
\begin{equation}
    \hat{\bf S}_{j} 
    = \sqrt{S_{j}} 
    \left[ 
    \Big( 1 - \frac{\hat{n}_{j}}{2S_{j}} \Big)^{1/2} \hat{a}_{j} \hat{\bf e}_{j}^{-} + 
    \hat{a}_{j}^{\dagger} \Big( 1 - \frac{\hat{n}_{j}}{2S_{j}} \Big)^{1/2} \hat{\bf e}_{j}^{+}
    \right]
    + (S_{j} - \hat{n}_{j}) \hat{\bf e}_{j}^{0},
\end{equation}
where the operators $\hat{a}_{i}$ and $\hat{a}_{i}^{\dagger}$ are the bosonic annihilation and creation operators that satisfy the canonical commutation relation $[\hat{a}_{i}, \hat{a}_{j}^{\dagger}] = \delta_{ij}$, and $\hat{n}_{j} \equiv \hat{a}^{\dagger}_{j} \hat{a}_{j}$ is the boson number operator. 
The vectors $\hat{\bf e}_{j}^{\alpha}$ ($\alpha = \pm, 0$) define a local frame of reference. In a more conventional Cartesian basis, one defines $\hat{\bf e}_{j}^{\pm} = (\hat{\bf x}_{j} \pm i \hat{\bf y}_{j})/\sqrt{2}$ and $\hat{\bf e}_{j}^{0} = \hat{\bf z}_{j}$ with $\hat{\bf e}_{j}^{-} \times \hat{\bf e}_{j}^{+} = i \hat{\bf e}_{j}^{0}$, where $\hat{\bf e}^{0}_{j} = {\bf S}_{j} / {S_{j}} = \hat{\bf n}_{j}$ denotes the classical spin orientation.

The perpendicular component $ \delta \hat{\mathbf{S}}_{j}^{\perp}$ satisfies $ \hat{\bf e}_{j}^{0} \cdot \delta \hat{\mathbf{S}}_{j}^{\perp} = 0 $ (where $\hat{\bf e}_{j}^{0}$ is the unit vector along the classical spin direction) and contains only odd-order bosonic operators: 
\begin{equation}
    \delta \hat{\mathbf{S}}_{j}^{\perp} = \sqrt{S_{j}}\left[ 
    \Big( 1 - \frac{\hat{n}_{j}}{2S_{j}} \Big)^{1/2} \hat{a}_{j} \hat{\bf e}_{j}^{-} + 
    \hat{a}_{j}^{\dagger} \Big( 1 - \frac{\hat{n}_{j}}{2S_{j}} \Big)^{1/2} \hat{\bf e}_{j}^{+}
    \right] 
    = \delta \hat{\mathbf{S}}_{j}^{(1)} + \delta \hat{\mathbf{S}}_{j}^{(3)} + \cdots 
\end{equation}
The parallel component, $\delta \hat{\mathbf{S}}_{j}^{\parallel} = \delta \hat{\mathbf{S}}_{j}^{(2)} = -\hat{n}_{j} \hat{\bf e}_{j}^{0}$, involves the number operator, where $\hat{\bf e}_{j}^{0}$ is the unit vector along the classical spin.

Based on the above equations, the Hamiltonian can then be rewritten as:
\begin{equation}
    \begin{split}
        H = & E_{\rm cl}
        + \sum_{i} \Big( \sum_{j} \mathbf{S}_{j} \cdot \mathbf{J}_{ji} - \mathbf{h}_{i} \Big) \cdot \delta \hat{\mathbf{S}}_{i}
        + \sum_{i,j} \delta \hat{\mathbf{S}}_{i} \cdot \mathbf{J}_{ij} \cdot \delta \hat{\mathbf{S}}_{j}, \\
        = & E_{\rm cl} + H_{1} + H_{2} + H_{3} + H_{4} + \cdots,
    \end{split}
\end{equation}
where the terms could be obtained by plugging $\delta \hat{\bf S}_{j}^{\perp} = \delta \hat{\mathbf{S}}_{j}^{(1)} + \delta \hat{\mathbf{S}}_{j}^{(3)} + \cdots$, and $\delta \hat{\bf S}_{j}^{\parallel} = \delta \hat{\bf S}_{j}^{(2)}$.
\begin{equation}
    \begin{split}
        E_{\rm cl} 
        = & \sum_{ij} \mathbf{S}_i \cdot \mathbf{J}_{ij} \mathbf{S}_j - \sum_{i} \mathbf{h}_{i} \cdot \mathbf{S}_{i}
        = \sum_{ij}  \widetilde{J}_{ij}^{00} S_{i} S_{j} - \sum_{i} S_{i}^{0} h_{i}, \\
        H_{1} = & \sum_{i} \Big( \sum_{j \in \mathsf{L}(i)} \mathbf{S}_{j} \cdot \mathbf{J}_{ji} - \mathbf{h}_{i} \Big) \cdot \delta \hat{\mathbf{S}}_{i}^{(1)} 
        = \sum_{i} \frac{\partial E_{\rm cl}}{\partial \mathbf{S}_{i}} \cdot \delta \hat{\mathbf{S}}_{i}^{(1)}, \\
        H_{2} = & \sum_{ij} \delta \hat{\mathbf{S}}_{i}^{(1)} \cdot \mathbf{J}_{ij} \cdot \delta \hat{\mathbf{S}}_{j}^{(1)} 
        + \sum_{i} \Big( \sum_{j\in \mathsf{L}(i)} \mathbf{S}_{j} \cdot \mathbf{J}_{ji} - \mathbf{h}_{i} \Big) \cdot \delta \hat{\mathbf{S}}_{i}^{(2)}, \\
        H_{3} = & \sum_{ij} \big( \delta \hat{\mathbf{S}}_{i}^{(1)} \cdot \mathbf{J}_{ij} \cdot \delta \hat{\mathbf{S}}_{j}^{(2)} + \delta \hat{\mathbf{S}}_{i}^{(2)} \cdot \mathbf{J}_{ij} \cdot \delta \hat{\mathbf{S}}_{j}^{(1)} \big), 
        \quad \text{and so on.} 
    \end{split}
\end{equation}
Note that the linear boson term vanishes when the spin configuration corresponds to a classical extremum:
\begin{equation}
    \sum_{i} \frac{\partial E_{\rm cl}}{\partial \mathbf{S}_{i}} \cdot \delta \hat{\mathbf{S}}_{i}^{(1)} = 0.
\end{equation}
To demonstrate this, consider the Lagrangian with a Lagrange multiplier enforcing the spin magnitude constraint:
\begin{equation}
    \mathcal{L} = E_{\rm cl} - \sum_{j} \frac{\lambda_{j}}{2} \big( |\mathbf{S}_{j}|^{2} - S_{j}^{2} \big).
\end{equation}
Taking the derivative with respect to $\mathbf{S}_{i}$, we get:
\begin{equation}
    \frac{\partial \mathcal{L}}{\partial \mathbf{S}_{i}} = \frac{\partial E_{\rm cl}}{\partial \mathbf{S}_{i}} - \lambda_{i} \mathbf{S}_{i} = 0,
\end{equation}
where $\frac{\partial E_{\rm cl}}{\partial \mathbf{S}_{i}} = \sum_{j} \mathbf{S}_{j} \cdot \mathbf{J}_{ji} - \mathbf{h}_{i}$. 
Since $\delta \hat{\mathbf{S}}_{i}^{\perp}$ is perpendicular to $\mathbf{S}_{i}$, it follows that:
\begin{equation}
    \boxed{\frac{\partial E_{\rm cl}}{\partial \mathbf{S}_{i}} \cdot \delta \hat{\mathbf{S}}_{i}^{\perp} = \lambda_{i} \mathbf{S}_{i} \cdot \delta \hat{\mathbf{S}}_{i}^{\perp} = 0, \text{ for all } i.} 
\end{equation}
Thus, the linear boson term, as well as any odd terms related to classical energy differentiation, vanishes when the spin configuration is a classical stable point.


\subparagraph{Comment ---}
Alternatively, \emph{Dyson-Maleev} transformation~\cite{PhysRev.102.1217, maleev1958scattering} could map spin operators onto bosonic creation and annihilation operators ($a_{i}^{\dagger}, a_{i}$) as well:
\begin{equation}\label{eq: spin-to-boson mapping}
        S^{+}_{i} = \sqrt{2S_{i}} \Big( 1 - \frac{a^{\dagger}_{i}a_{i}}{2S_{i}} \Big)^{\xi} a_{i}, \quad
        S^{-}_{i} = \sqrt{2S_{i}} \, a_{i}^{\dagger} \Big( 1 - \frac{a^{\dagger}_{i}a_{i}}{2S_{i}} \Big)^{1 - \xi}, \quad
        S_{i}^{z} = S_{i} - a^{\dagger}_{i} a_{i},
\end{equation}
where $\xi = 1/2$ for the Holstein-Primakoff transformation, and $\xi = 1$ for the Dyson-Maleev transformation. 
While both transformations map spin operators onto bosonic operators, the Dyson-Maleev transformation does not preserve Hermiticity. For instance, $(S^{+}_{i})^{\dagger}$ is not equal to $S^{-}_{i}$:
\begin{equation}
        (S^{+}_{i})^{\dagger} = \sqrt{2S_{i}} a^{\dagger}_{i} \Big( 1 - \frac{a^{\dagger}_{i} a_{i}}{2S_{i}} \Big)^{\xi} = S^{-}_{i} \Big( 1 - \frac{a^{\dagger}_{i} a_{i}}{2S_{i}} \Big)^{2\xi - 1},
\end{equation}
whereas for $\xi = 1/2$ (the Holstein-Primakoff transformation), Hermiticity of the spin operators is preserved.




\subsection{Rotations for Spin Models}
To analyze an ordered spin system using linear spin-wave theory (LSWT), it is useful to rotate the spins at each site from the \emph{laboratory} reference frame $\{\hat{\bf x}, \hat{\bf y}, \hat{\bf z}\}$ to a \emph{local} reference frame $\{\hat{\bf x}_{j}, \hat{\bf y}_{j}, \hat{\bf z}_{j}\}$ (aligned with the classical spin ordering, where $\hat{{\bf z}}_{j}$ coincides with the spin's quantization axis). 
The transformation is given by:
\begin{equation}
    \label{eq: Spin rotation}
    \hat{\bf S}_j = \mathbf{R}_j \, \widetilde{\mathbf{S}}_j 
    = \widetilde{S}^{x}_{j} \hat{\bf x}_{j} + \widetilde{S}^{y}_{j} \hat{\bf y}_{j} + \widetilde{S}^{z}_{j} \hat{\bf z}_{j},
\end{equation}
where $\widetilde{\mathbf{S}}_j$ is the spin operator in the local reference frame at site $j$.
By applying the rotation to all spins in a unit cell, the spin Hamiltonian is transformed from the laboratory frame to the local frame as:
\begin{equation}\label{eq: Rotated Hamiltonian}
    \begin{split}
        \hat{H} = & \sum_{ij} \sum_{\alpha, \beta} S^{\alpha}_{i} J_{ij}^{\alpha\beta} S^{\beta}_{j} - \sum_{j, \alpha} h_{j}^{\alpha} S^{\alpha}_{j}
        = \sum_{ij} \sum_{\alpha, \beta} \widetilde{S}^{\alpha}_{i} \widetilde{J}_{ij}^{\alpha\beta} \widetilde{S}^{\beta}_{j} - \sum_{j, \alpha} \widetilde{h}_{j}^{\alpha} \widetilde{S}^{\alpha}_{j},
    \end{split}
\end{equation}
where the rotated couplings are defined as:
\begin{equation}
    \widetilde{J}_{ij}^{\alpha\beta}  = [\widetilde{\mathbf{J}}_{ij}]^{\alpha\beta} = [\mathbf{R}_i^{\rm T} \mathbf{J}_{ij} \mathbf{R}_j]^{\alpha\beta}, 
    \quad \text{ and } \quad
     \widetilde{h}_{j}^{\alpha} = [\widetilde{\mathbf{h}}_j]^{\alpha} = [\mathbf{h}_{j}^{\rm T} \mathbf{R}_{j}]^{\alpha}.
\end{equation}
In practice, the rotation from the crystal frame to the local frame is performed in two stages, expressed as $\mathbf{R}_i = \mathbf{R}^{\varphi} \cdot \mathbf{R}^{\theta}$. The combined rotation matrix is:
\begin{equation}
    \mathbf{R} (\theta, \varphi) = \mathbf{R}^{\varphi} \cdot \mathbf{R}^{\theta} = 
    \begin{pmatrix} 
        \cos \theta \cos \varphi & -\sin \varphi & \sin \theta \cos \varphi \\
        \cos \theta \sin \varphi & \cos \varphi & \sin \theta \sin \varphi \\
        -\sin \theta & 0 & \cos \theta
    \end{pmatrix},
    \label{eq_frametransform_gen}
\end{equation}
where $\theta$ is the angle relative to the $z$-axis in lab frame and $\varphi$ is the azimuthal angle in the $xy$-plane.
Note that the classical spin orientation $\mathbf{S}_j$ satisfies:
\begin{equation}
    \mathbf{S}_{j} 
    = \mathbf{R}_{j} \, \hat{\bf z} 
    = \mathbf{R}^{\varphi_{j}} \mathbf{R}^{\theta_{j}} 
    \begin{pmatrix} 
        0 \\ 0 \\ 1 
    \end{pmatrix} 
    = \mathbf{R}^{\varphi_j} 
    \begin{pmatrix} 
        \sin \theta_{j} \\ 0 \\ \cos \theta_{j} 
    \end{pmatrix} 
    = 
    \begin{pmatrix} 
        \sin \theta_{j} \cos \varphi_{j} \\ \sin \theta_{j} \sin \varphi_{j} \\ \cos \theta_{j} 
    \end{pmatrix}.
\end{equation}




For later uses, it is also convenient to express the spin operator using ladder operators with complex vector:
\begin{equation}\label{eq: Spin rotation-ladder}
    \hat{\bf S}_j = \mathbf{R}_j \, \widetilde{\mathbf{S}}_j
    = \frac{\widetilde{S}_{j}^{+} \hat{\bf e}^{-}_{j} + \widetilde{S}_{j}^{-} \hat{\bf e}^{+}_{j} }{\sqrt{2}} + \widetilde{S}^{0}_{j} \hat{\bf e}^{0}_{j},
\end{equation}
where the spin operators are $\widetilde{S}_{j}^{\pm} \equiv \widetilde{S}_{j}^{x} \pm i \widetilde{S}_{j}^{y}$, $\widetilde{S}^{0}_{j} = \widetilde{S}^{z}_{j}$, and $[\widetilde{S}^{z}_{j}, \widetilde{S}^{\pm}_{j}] = \pm \widetilde{S}^{\pm}_{j}$.
The complex vectors are
\begin{equation}
    \hat{\bf e}_{j}^{\alpha} \equiv \mathbf{R}_{j} \hat{\bf e}^{\alpha}: \quad 
    \hat{\bf e}^{\pm}_{j} = \mathbf{R}_{j} \Big( \frac{ \hat{\bf x} \pm i \hat{\bf y} }{\sqrt{2}} \Big)
    \text{ and }
    \hat{\bf e}_{j}^{0} \equiv \mathbf{R}_{j} \hat{\bf z},
\end{equation}
where $\alpha = \pm, 0, (x,y,z)$ and the $\hat{\bf e}_{j}^{0}$-direction represents the classical magnetic orientation at site $j$. 
These complex vectors $\hat{\bf e}_{j}^{\alpha}$ introduce a rotated Hamiltonian similar to the previous one:
\begin{equation}
    \hat{H} 
    = \sum_{ij} \sum_{\alpha, \beta = 0, \pm} 
    \widetilde{S}^{\alpha}_{i} \widetilde{J}^{\bar{\alpha}\bar{\beta}}_{ij} \widetilde{S}^{\beta}_{j} 
    - \sum_{j} \sum_{\alpha = 0, \pm} \widetilde{h}^{\alpha}_{j} \widetilde{S}^{\alpha}_{j}.
\end{equation}
where $\bar{(\cdot)} = - ( \cdot )$, for example $\alpha = \pm, (0)$ and $\bar{\alpha} = \mp, (0)$.
The couplings $\widetilde{J}^{\bar{\alpha}\bar{\beta}}_{ij}$ and magnetic fields $\widetilde{h}_{j}^{\alpha}$ are defined as 
\begin{equation}
    \widetilde{J}^{\alpha\beta}_{ij} \equiv \hat{\bf e}_{i}^{\alpha} \cdot \widetilde{\bf J}_{ij} \cdot \hat{\bf e}_{j}^{\beta}
    \quad \text{ and } \quad
    \widetilde{h}^{\alpha}_{j} \equiv \widetilde{\bf h}_{j} \cdot \hat{\bf e}^{\alpha}_{j}.
\end{equation}

The transformation from $xyz$ to $\pm0$ basis is given by:
\begin{equation}
    \mathbf{C} = 
    \begin{pmatrix}
        \hat{\bf e}^{+} & \hat{\bf e}^{-} & \hat{\bf e}^{0}
    \end{pmatrix}
    = 
    \begin{pmatrix}
        [\hat{\mathbf{e}}^{+}]^{x} & [\hat{\mathbf{e}}^{-}]^{x} & [\hat{\mathbf{e}}^{0}]^{x} \\[0.1em]
        [\hat{\mathbf{e}}^{+}]^{y} & [\hat{\mathbf{e}}^{-}]^{y} & [\hat{\mathbf{e}}^{0}]^{y} \\[0.1em]
        [\hat{\mathbf{e}}^{+}]^{z} & [\hat{\mathbf{e}}^{-}]^{z} & [\hat{\mathbf{e}}^{0}]^{z} 
    \end{pmatrix}
    = 
    \begin{pmatrix}
        \frac{1}{\sqrt{2}}  &   \frac{1}{\sqrt{2}}  &   0   \\[0.1em]
        \frac{i}{\sqrt{2}}  & - \frac{i}{\sqrt{2}}  &   0   \\[0.1em]
                    0       &           0           &   1   
    \end{pmatrix} 
\end{equation}
where $[\hat{\mathbf{e}}_{i}^{\alpha}]^{\beta}$ denotes the $\beta = x, y, z$ component of the basis vector $\hat{\mathbf{e}}_{i}^{\alpha}$ (with $\alpha = +, -, 0$). 
These components are obtained from:
\begin{equation}
    \mathbf{U}_{j} = \mathbf{R}_{j} \mathbf{C} =
    \begin{pmatrix}
        [\hat{\mathbf{e}}_{j}^{+}]^{x} & [\hat{\mathbf{e}}_{j}^{-}]^{x} & [\hat{\mathbf{e}}_{j}^{0}]^{x} \\[0.25em]
        [\hat{\mathbf{e}}_{j}^{+}]^{y} & [\hat{\mathbf{e}}_{j}^{-}]^{y} & [\hat{\mathbf{e}}_{j}^{0}]^{y} \\[0.25em]
        [\hat{\mathbf{e}}_{j}^{+}]^{z} & [\hat{\mathbf{e}}_{j}^{-}]^{z} & [\hat{\mathbf{e}}_{j}^{0}]^{z} 
    \end{pmatrix}
    = 
    \begin{pmatrix}
        R^{xx}_{j}   &   R^{xy}_{j}   &   R^{xz}_{j}   \\[0.25em]
        R^{yx}_{j}   &   R^{yy}_{j}   &   R^{yz}_{j}   \\[0.25em]
        R^{zx}_{j}   &   R^{zy}_{j}   &   R^{zz}_{j}   
    \end{pmatrix}
    \begin{pmatrix}
        \frac{1}{\sqrt{2}}   &   \frac{1}{\sqrt{2}}   &   0   \\[0.25em]
        \frac{i}{\sqrt{2}}   & - \frac{i}{\sqrt{2}}   &   0   \\[0.25em]
                0     &         0       &   1   
    \end{pmatrix} 
\end{equation}

The transformed coupling tensor $\mathbf{C}^{\dagger} \mathbf{J} \mathbf{C} = $ is expressed as:
\begin{equation}
    \begin{split}
    \begin{pmatrix}
        J^{-+}_{ij} & J^{--}_{ij} & J^{-0}_{ij} \\[0.3em]
        J^{++}_{ij} & J^{+-}_{ij} & J^{+0}_{ij} \\[0.3em]
        J^{0+}_{ij} & J^{0-}_{ij} & J^{00}_{ij}
    \end{pmatrix}
    = &
    \begin{pmatrix}
        \frac{1}{\sqrt{2}}   & - \frac{i}{\sqrt{2}}   &   0   \\[0.3em]
        \frac{1}{\sqrt{2}}   &   \frac{i}{\sqrt{2}}   &   0   \\[0.3em]
                0     &         0       &   1   
    \end{pmatrix} 
    \begin{pmatrix}
        J^{xx}_{ij}   &   J^{xy}_{ij}   &   J^{xz}_{ij}   \\[0.3em]
        J^{yx}_{ij}   &   J^{yy}_{ij}   &   J^{yz}_{ij}   \\[0.3em]
        J^{zx}_{ij}   &   J^{zy}_{ij}   &   J^{zz}_{ij}   
    \end{pmatrix}
    \begin{pmatrix}
        \frac{1}{\sqrt{2}}   &   \frac{1}{\sqrt{2}}   &   0   \\[0.3em]
        \frac{i}{\sqrt{2}}   & - \frac{i}{\sqrt{2}}   &   0   \\[0.3em]
                0     &         0       &   1   
    \end{pmatrix} 
    % \\ 
    % \begin{pmatrix}
    %     \frac{1}{2} \left( J^{xx} + J^{yy} + i (J^{xy} - J^{yx}) \right) & 
    %     \frac{1}{2} \left( J^{xx} - J^{yy} - i (J^{xy} + J^{yx}) \right) & 
    %     \frac{1}{\sqrt{2}} \left( J^{xz} - i J^{yz} \right) \\
    %     \frac{1}{2} \left( J^{xx} - J^{yy} + i (J^{xy} + J^{yx}) \right) & 
    %     \frac{1}{2} \left( J^{xx} + J^{yy} - i (J^{xy} - J^{yx}) \right) & 
    %     \frac{1}{\sqrt{2}} \left( J^{xz} + i J^{yz} \right) \\
    %     \frac{1}{\sqrt{2}} \left( J^{zx} + i J^{zy} \right) & 
    %     \frac{1}{\sqrt{2}} \left( J^{zx} - i J^{zy} \right) & 
    %     J^{zz}
    % \end{pmatrix}
    \end{split}
\end{equation}
with elements defined as $\widetilde{J}_{ij}^{\alpha\beta} = \hat{\mathbf{e}}_i^{\alpha} \cdot \widetilde{\mathbf{J}}_{ij} \cdot \hat{\mathbf{e}}_j^{\beta}$:
\begin{equation}\label{eq: expressions for exchange couplings}
    \begin{array}{rlrlrl}
        J_{ij}^{--} &= \frac{1}{2} \big( J_{ij}^{xx} - J_{ij}^{yy} - i (J_{ij}^{xy} + J_{ij}^{yx}) \big), \qquad \qquad &
        J_{ij}^{++} &= \frac{1}{2} \big( J_{ij}^{xx} - J_{ij}^{yy} + i (J_{ij}^{xy} + J_{ij}^{yx}) \big), &&\\[0.25em]
        J_{ij}^{-+} &= \frac{1}{2} \big( J_{ij}^{xx} + J_{ij}^{yy} + i (J_{ij}^{xy} - J_{ij}^{yx}) \big), \qquad &
        J_{ij}^{+-} &= \frac{1}{2} \big( J_{ij}^{xx} + J_{ij}^{yy} - i (J_{ij}^{xy} - J_{ij}^{yx}) \big), &&\\[0.25em]
        J_{ij}^{0\pm} &= \frac{1}{\sqrt{2}} \big( J_{ij}^{zx} \pm i J_{ij}^{zy} \big), &
        J_{ij}^{\pm0} &= \frac{1}{\sqrt{2}} \big( J_{ij}^{xz} \pm i J_{ij}^{yz} \big), & \qquad
        J_{ij}^{00} &= J_{ij}^{zz}.
    \end{array}
\end{equation}





\subsection{Bosonic representations for spin Hamiltonian}\label{subsec: Bosonic representations for spin exchange Hamiltonian}

In this section, we derive the explicit bosonic Hamiltonian for the spin model.
\begin{equation}
    \begin{split}
        \hat{H} 
        &= E_{\rm cl}
        + \sum_{i} \Big( \sum_{j \in \mathsf{L}(ij)} \mathbf{S}_{j} \cdot \mathbf{J}_{ji} - \mathbf{h}_{i} \Big) \cdot \delta \hat{\mathbf{S}}_{i}
        + \sum_{ij \in \mathsf{L}} \delta \hat{\mathbf{S}}_{i} \cdot \mathbf{J}_{ij} \cdot \delta \hat{\mathbf{S}}_{j},
    \end{split}
\end{equation}
where $E_{\rm cl}$ is the classical energy, $\mathbf{S}_{j}$ represents the spin vector at site $j$, $\mathbf{J}_{ji}$ is the exchange interaction tensor, $\mathbf{h}_{i}$ is the external field, and $\delta \hat{\mathbf{S}}_{i}$ denotes the spin fluctuation operator.

Let us divide the spin model into even and odd bosonic operator terms,
\begin{subequations}\label{eq: LSWT boson expansion even/odd}
    \begin{align}
        H_{\rm even} &= \sum_{ij\in\mathsf{L}} 
        \Big( \delta \hat{\mathbf{S}}_{i}^{\perp} \cdot \mathbf{J}_{ij} \cdot \delta \hat{\mathbf{S}}_{j}^{\perp} 
        + \delta \hat{\mathbf{S}}_{i}^{\parallel} \cdot \mathbf{J}_{ij} \cdot \delta \hat{\mathbf{S}}_{j}^{\parallel} \Big)
        + \sum_{i} \Big( \sum_{j \in \mathsf{L}(ij)} \mathbf{S}_{j} \cdot \mathbf{J}_{ji} - \mathbf{h}_{i} \Big) \cdot \delta \hat{\mathbf{S}}_{i}^{\parallel}, 
        \label{subeq: even bosonic terms in SWT} \\
        H_{\rm odd}^{(m \ge 3)} &= \sum_{ij\in\mathsf{L}} 
        \Big( \delta \hat{\mathbf{S}}_{i}^{\perp} \cdot \mathbf{J}_{ij} \cdot \delta \hat{\mathbf{S}}_{j}^{\parallel} 
        + \delta \hat{\mathbf{S}}_{i}^{\parallel} \cdot \mathbf{J}_{ij} \cdot \delta \hat{\mathbf{S}}_{j}^{\perp} \Big).
        \label{subeq: odd bosonic terms in SWT} 
    \end{align}
\end{subequations}
and substituting the following spin-to-boson transformation:
\begin{equation}
    \delta \hat{\mathbf{S}}_{i}^{\perp} = \frac{\widetilde{S}_{i}^{+} \hat{\mathbf{e}}^{-}_{i} + \widetilde{S}_{i}^{-} \hat{\mathbf{e}}^{+}_{i} }{\sqrt{2}} = \sqrt{S_{i}} \left[ \Big( 1 - \frac{\hat{a}^{\dagger}_{i}\hat{a}_{i}}{2S_{i}} \Big)^{1/2} \hat{a}_{i} \hat{\mathbf{e}}_{i}^{-} + \hat{a}_{i}^{\dagger} \Big( 1 - \frac{\hat{a}^{\dagger}_{i}\hat{a}_{i}}{2S_{i}} \Big)^{1/2} \hat{\mathbf{e}}_{i}^{+} \right], 
    \quad \text{ and } \quad
    \delta \hat{\mathbf{S}}^{\parallel}_{i} =  - \hat{n}_{i} \hat{\mathbf{e}}_{i},
\end{equation}
where $\hat{\mathbf{e}}_{i}^{\pm}$ are the raising and lowering unit vectors, $\hat{\mathbf{e}}_{i}$ is the unit vector along the local quantization axis, and $\hat{n}_{i} = \hat{a}^{\dagger}_{i}\hat{a}_{i}$ is the boson number operator.
This boson mapping results in the SWT Hamiltonian. 
While the odd terms can also be obtained easily, we will focus on the even terms as they are of primary interest for the linear spin wave theory.

The even Hamiltonian term, $H_{\rm even}$, consists of terms involving an even number of bosonic operators: 
\begin{equation}
    \delta \hat{\mathbf{S}}_{i}^{\perp} \cdot \mathbf{J}_{ij} \cdot \delta \hat{\mathbf{S}}_{j}^{\perp}  
    = \frac{1}{2} \Big( \widetilde{S}_{i}^{+} \widetilde{J}_{ij}^{--} \widetilde{S}_{j}^{+}  
    + \widetilde{S}_{i}^{+} \widetilde{J}_{ij}^{-+} \widetilde{S}_{j}^{-}
    + \widetilde{S}_{i}^{-} \widetilde{J}_{ij}^{+-} \widetilde{S}_{j}^{+} 
    + \widetilde{S}_{i}^{-} \widetilde{J}_{ij}^{++} \widetilde{S}_{j}^{-} \Big),
\end{equation}
This becomes:
\begin{align*}
    \frac{1}{2} \widetilde{S}_{i}^{+} \widetilde{J}_{ij}^{-+} \widetilde{S}_{j}^{-} 
    = & 
    \sqrt{S_{i}S_{j}} \widetilde{J}_{ij}^{-+} \Big( 1 - \frac{\hat{n}_{i}}{2S_i} \Big)^{1/2} \hat{a}_{i} \hat{a}_{j}^{\dagger}  \Big( 1 - \frac{\hat{n}_{j}}{2S_j} \Big)^{1/2} 
    = t_{ij}^{-+} \Big[ \hat{a}_{i} \hat{a}_{j}^{\dagger}  - \tfrac{1}{4} \Big( \tfrac{\hat{n}_{i} \hat{a}_{i} \hat{a}_{j}^{\dagger}}{S_{i}} + \tfrac{\hat{a}_{i} \hat{a}_{j}^{\dagger} \hat{n}_{j}}{S_{j}} \Big) + \mathcal{O} (S^{-2}) \Big]
    \\
    \frac{1}{2} \widetilde{S}_{i}^{-} \widetilde{J}_{ij}^{+-} \widetilde{S}_{j}^{+} = &
    \sqrt{S_i S_j} \widetilde{J}_{ij}^{+-} \hat{a}_i^{\dagger} \Big( 1 - \frac{\hat{n}_i}{2 S_i} \Big)^{1/2} \Big( 1 - \frac{\hat{n}_j}{2 S_j} \Big)^{1/2} \hat{a}_j 
    = t_{ij}^{+-} \Big[ \hat{a}_{i}^{\dagger} \hat{a}_{j}  - \tfrac{1}{4} \Big( \tfrac{\hat{a}_{i}^{\dagger} \hat{n}_{i} \hat{a}_{j}}{S_{i}} + \tfrac{\hat{a}_{i}^{\dagger} \hat{n}_{j} \hat{a}_{j}}{S_{j}} \Big) + \mathcal{O} (S^{-2}) \Big]
    \\
    \frac{1}{2} \widetilde{S}_{i}^{+} \widetilde{J}_{ij}^{--} \widetilde{S}_{j}^{+} = & 
    \sqrt{S_i S_j} \widetilde{J}_{ij}^{--} \Big( 1 - \frac{\hat{n}_i}{2 S_i} \Big)^{1/2} \hat{a}_i \Big( 1 - \frac{\hat{n}_j}{2 S_j} \Big)^{1/2} \hat{a}_j 
    = t_{ij}^{--} \Big[ \hat{a}_{i} \hat{a}_{j} - \tfrac{1}{4} \Big( \tfrac{\hat{n}_{i} \hat{a}_{i} \hat{a}_{j}}{S_{i}} + \tfrac{\hat{a}_{i} \hat{n}_{j} \hat{a}_{j}}{S_{j}} \Big) + \mathcal{O} (S^{-2}) \Big]
    \\
    \frac{1}{2} \widetilde{S}_{i}^{-} \widetilde{J}_{ij}^{++} \widetilde{S}_{j}^{-} = & 
    \sqrt{S_i S_j} \widetilde{J}_{ij}^{++} \hat{a}_i^{\dagger} \Big( 1 - \frac{\hat{n}_i}{2 S_i} \Big)^{1/2} \hat{a}_j^{\dagger} \Big( 1 - \frac{\hat{n}_j}{2 S_j} \Big)^{1/2}
    = t_{ij}^{++} \Big[ \hat{a}_{i}^{\dagger} \hat{a}_{j}^{\dagger} - \tfrac{1}{4} \Big( \tfrac{\hat{a}_{i}^{\dagger} \hat{n}_{i} \hat{a}_{j}^{\dagger}}{S_{i}} + \tfrac{\hat{a}_{i}^{\dagger} \hat{a}_{j}^{\dagger} \hat{n}_{j}}{S_{j}} \Big) + \mathcal{O} (S^{-2}) \Big],
\end{align*}
where we define $t^{\alpha\beta}_{ij} = \sqrt{S_{i}S_{j}} \widetilde{J}^{\alpha\beta}_{ij}$.
Next, the other terms in Eqn.~\eqref{subeq: even bosonic terms in SWT} are:
\begin{equation}
    \begin{split}
    \sum_{i} \Big( \sum_{j \in \mathsf{L}(i)} \mathbf{S}_{j} \cdot \mathbf{J}_{ji} - \mathbf{h}_{i} \Big) \cdot \delta \hat{\mathbf{S}}_{i}^{\parallel} 
    &= 
    \sum_{i} \Big( \widetilde{h}_{i}^{0} - \sum_{j \in \mathsf{L}(i)} S_{j} \cdot \widetilde{J}_{ji}^{00} \Big)  \hat{n}_{i}, \\
    \delta \hat{\mathbf{S}}_{i}^{\parallel} \cdot \mathbf{J}_{ij} \cdot \delta \hat{\mathbf{S}}_{j}^{\parallel}
    &= \hat{n}_{i} \widetilde{J}_{ij}^{00} \hat{n}_{j},
    \end{split}
\end{equation}
where $\widetilde{h}_{i}^{0}$ is the magnetic field along the local quantization axis and $\widetilde{J}_{ij}^{00}$ represents the longitudinal component of the exchange interaction.

Combining all terms, we can express the quadratic and quartic Hamiltonian contributions, $H_{2}$ and $H_{4}$, as:
\begin{align}
    H_{2} &= \sum_{ij} \Big[ t_{ij}^{--} \hat{a}_{i} \hat{a}_{j} + t_{ij}^{-+} \hat{a}_{i}\hat{a}_{j}^{\dagger} + t_{ij}^{+-} \hat{a}_{i}^{\dagger} \hat{a}_{j} + t_{ij}^{++} \hat{a}^{\dagger}_{i} \hat{a}^{\dagger}_{j} - \widetilde{J}_{ij}^{00} \big( S_{i} \hat{n}_{j} + S_{j} \hat{n}_{i} \big) \Big] + \sum_{i} h_{i} \hat{n}_{i}, 
    \label{eq: SWT bosonic H2 in real space}\\
    H_{4} &= \sum_{ij} \widetilde{J}_{ij}^{00} \hat{n}_{i} \hat{n}_{j} - \Big[ t_{ij}^{--} \big( \tfrac{\hat{n}_{i}}{S_{i}} + \tfrac{\hat{n}_{j}}{S_{j}} \big) \hat{a}_{i} \hat{a}_{j} + t_{ij}^{++} \hat{a}^{\dagger}_{i} \hat{a}^{\dagger}_{j} \big( \tfrac{\hat{n}_{i}}{S_{i}} + \tfrac{\hat{n}_{j}}{S_{j}} \big) + t_{ij}^{-+} \big( \tfrac{\hat{n}_{i} \hat{a}_{i} \hat{a}_{j}^{\dagger}}{S_{i}} + \tfrac{\hat{a}_{i} \hat{a}_{j}^{\dagger} \hat{n}_{j}}{S_{j}} \big) + t_{ij}^{+-} \big( \tfrac{\hat{a}_{i}^{\dagger} \hat{n}_{i} \hat{a}_{j}}{S_{i}} + \tfrac{\hat{a}_{i}^{\dagger} \hat{n}_{j} \hat{a}_{j}}{S_{j}} \big) \Big]. 
    \label{eq: SWT bosonic H4 in real space}
\end{align}
Note that if the anisotropic couplings are absent in the rotated exchange matrix $\widetilde{J}^{xy} = \widetilde{J}^{yx} = 0$, and $\widetilde{J}^{xx} = \widetilde{J}^{yy}$ are equal, then the anomalous contribution, such as $a^{\dagger}a^{\dagger}$, vanishes, see Eq.~\eqref{eq: expressions for exchange couplings}. 
Furthermore, when the spins' magnitudes equal $S_{i} = S$, the expansions 
\begin{equation}
    H = S^{2} \tilde{E}_{cl} + S^{1} \tilde{H}_2 + S^{1/2} \tilde{H}_3 + S^{0} \tilde{H}_4 + \cdots.
\end{equation}


\subsection{Momentum Space Representations}

In this section, we shall find the momentum space representation for the quadratic spin wave Hamiltonian ($H_{2}$).
The quadratic bosonic Hamiltonian we have in Eqn.~\eqref{eq: SWT bosonic H2 in real space} is 
\begin{equation}
    \begin{split}
        H_{2} = & \sum_{ij}
        \Big[ t_{ij}^{--} \hat{a}_{i} \hat{a}_{j} + t_{ij}^{-+} \hat{a}_{i}\hat{a}_{j}^{\dagger} + t_{ij}^{+-} \hat{a}_{i}^{\dagger} \hat{a}_{j} + t_{ij}^{++} \hat{a}^{\dagger}_{i} \hat{a}^{\dagger}_{j} - \widetilde{J}_{ij}^{00} \big( S_{i} \hat{a}_{j}^{\dagger} \hat{a}_{j} + S_{j} \hat{a}_{i}^{\dagger} \hat{a}_{i} \big) \Big]  + \sum_{i} \widetilde{h}^{0}_{i} \hat{a}^{\dagger}_{i}\hat{a}_{i} \\
        = & \sum_{ij} 
        \Big[ t_{ij}^{--} \hat{a}_{i} \hat{a}_{j} + t_{ij}^{-+} \hat{a}_{i}\hat{a}_{j}^{\dagger} + t_{ij}^{+-} \hat{a}_{i}^{\dagger} \hat{a}_{j} + t_{ij}^{++} \hat{a}^{\dagger}_{i} \hat{a}^{\dagger}_{j} \Big] 
        + \sum_{i} \mu_{i} \hat{a}^{\dagger}_{i} \hat{a}_{i}, 
    \end{split}
\end{equation}
where the (effective) chemical potential $\mu_{i}$ is 
\begin{equation}
    \mu_{i} = \widetilde{h}^{0}_{i} - \sum_{j \in \mathsf{L}(i)} S_{j} \widetilde{J}^{00}_{ij} = \frac{\partial E_{\rm cl}}{\partial \mathbf{S}_{i}} \cdot \hat{\mathbf{e}}_{i}^{0}.
\end{equation}

We assume translational invariance.  
In general, the spin system has magnetic sublattices $\sigma = a, b, c, \dots$.
The bosonic annihilation operator $\hat{a}_{i}$ is defined with two indices $i = (i, \sigma)$, where $\sigma$ denotes the magnetic sublattice and $i$ is the unit-cell index. For example, $\hat{b}_{(i, a)} = \hat{a}_{i}$ and $\hat{b}_{(i, b)} = \hat{b}_{i}$. 
This allows the operators to be transformed into their momentum-space representations:
\begin{equation}
    \hat{b}_{i\sigma} = \frac{1}{\sqrt{L}} \sum_{\mathbf{k} \in \mathrm{FBZ}} e^{i \mathbf{k} \cdot \mathbf{r}_{i\sigma}} \hat{b}_{\mathbf{k}\sigma},
\end{equation}
where $L$ is the number of unit cells in the system and FBZ denotes the first Brillouin zone.
Then, the Hamiltonian becomes:
\begin{align}
    H_{2} 
    = & \sum_{\mathbf{k}} 
    \Big[ \sum_{\mathsf{L}(i\sigma, j\rho)}  
        t^{+-}_{ij} e^{-i \mathbf{k} \cdot \bm{\delta}_{ij}} \hat{a}_{\mathbf{k} \sigma}^{\dagger}          \hat{a}_{ \mathbf{k} \rho }^{\phantom \dagger}   +
        t^{-+}_{ij} e^{ i \mathbf{k} \cdot \bm{\delta}_{ij}} \hat{a}_{\mathbf{k} \sigma}^{\phantom \dagger} \hat{a}_{ \mathbf{k} \rho }^{\dagger}  +
        t^{++}_{ij} e^{-i \mathbf{k} \cdot \bm{\delta}_{ij}} \hat{a}_{\mathbf{k} \sigma}^{\dagger}          \hat{a}_{-\mathbf{k} \rho }^{\dagger} + t^{--}_{ij} e^{ i \mathbf{k} \cdot \bm{\delta}_{ij}} \hat{a}_{\mathbf{k} \sigma}^{\phantom \dagger} \hat{a}_{-\mathbf{k} \rho }^{\phantom \dagger} \Big] 
    + \sum_{\mathbf{k}, \sigma} \mu_{\sigma} \hat{a}_{\mathbf{k}\sigma}^{\dagger} \hat{a}_{\mathbf{k}\sigma}^{\phantom\dagger} \nonumber \\
    = & \frac{1}{2} \sum_{\mathbf{k} \in \mathrm{FBZ}} 
    \begin{pmatrix}
        \psi_{  \mathbf{k}}^{\dagger} &  
        \psi_{-\mathbf{k}}
    \end{pmatrix}
    \begin{pmatrix}
        \mathsf{A}_{\mathbf{k}}          & \mathsf{B}_{\mathbf{k}} \\
        \mathsf{B}_{\mathbf{k}}^{\dagger}& \mathsf{A}_{-\mathbf{k}}^{*}
    \end{pmatrix}
    \begin{pmatrix}
        \psi_{  \mathbf{k}} \\
        \psi_{-\mathbf{k}}^{\dagger} 
    \end{pmatrix}
    + \frac{1}{2} \sum_{\mathbf{k} \in \mathrm{FBZ}}  C_{\mathbf{k}},
\end{align}
where $\psi^{\dagger}_{\mathbf{k}} = (a^{\dagger}_{\mathbf{k}}, b^{\dagger}_{\mathbf{k}}, \cdots)$ is a vector of bosonic creation operators for each sublattice, and the constant $C_{\mathbf{k}}$ follows from the normal ordering $\hat{a}_{\mathbf{k}\rho}^{\dagger} \hat{a}_{\mathbf{k}\sigma} = \frac{1}{2} (\hat{a}_{\mathbf{k}\rho}^{\dagger} \hat{a}_{ \mathbf{k}\sigma} + \hat{a}_{ \mathbf{k}\rho} \hat{a}_{\mathbf{k}\sigma}^{\dagger} ) - \frac{1}{2} \delta_{\rho\sigma}$.
\begin{align}
    \sum_{\mathbf{k} \in \mathrm{FBZ}} C_{\mathbf{k}} 
    = & \sum_{\mathbf{k} \in \mathrm{FBZ}} \Big[
    \sum_{\mathsf{L}(i\sigma,j\rho)} \big( t^{-+}_{ij} e^{i \mathbf{k} \cdot \bm{\delta}_{ij}} - t^{+-}_{ij} e^{-i \mathbf{k} \cdot \bm{\delta}_{ij}}\big) \delta_{\rho\sigma}    - \sum_{\sigma} \mu_{\sigma} \Big]
    = - L \sum_{\sigma = a, b, \cdots} \mu_{\sigma} \nonumber \\
    = & L \sum_{\sigma = a, b, \cdots} \Big( \sum_{\rho \in \mathsf{L}(\sigma)} S_{\rho} \widetilde{J}_{\rho\sigma}^{00} - \widetilde{h}_{\sigma}^{0} \Big) 
    = \sum_{i} \frac{\partial E_{\rm cl}}{\partial \mathbf{S}_{i}} \cdot \hat{\mathbf{e}}_{i}^{0} 
\end{align}
In the first line, the first term always vanishes if $\rho \neq \sigma$.
Even when $\rho = \sigma$ for hopping within the same sublattice but between different unit cells, these terms vanish in the summation over the FBZ.
Note that the constant term can be expressed as $\frac{1}{2} \sum_{\mathbf{k} \in \mathrm{FBZ}}  C_{\mathbf{k}}  = - \frac{1}{4} \Tr[\mathsf{H}_{\mathbf{k}}]$ and 
\begin{equation}
    \frac{1}{2} \sum_{\mathbf{k} \in \mathrm{FBZ}}  C_{\mathbf{k}} 
    = \frac{1}{2} \sum_{i} \frac{\partial E_{\rm cl}}{\partial \mathbf{S}_{i}} \cdot \hat{\mathbf{e}}_{i}^{0} 
    = \frac{1}{2} \sum_{i} \Big( \widetilde{h}_{i}^{0} - \sum_{j \in \mathsf{L}(i)} S_{j} \cdot \widetilde{J}_{ji}^{00} \Big)
    \xrightarrow{ S_{i} = S}  
    \frac{1}{S} \big( E_{\rm exc} + \frac{1}{2}  E_{\rm zm} \big),
\end{equation}
where $ E_{\rm exc} = \sum_{ij} \mathbf{S}_i \cdot \mathbf{J}_{ij} \cdot \mathbf{S}_j $ is the exchange energy and $ E_{\rm zm} = - \sum_{i} \mathbf{h}_{i} \cdot \mathbf{S}_{i}$ is the Zeeman energy. 
This is the reason why the quantum correction reduces the classical energy. 
Clearly, when $\mathbf{h} = 0$, we have
\begin{equation}
    \hat{H} = S(S+1) E_{\rm cl} + S H_{2} + S^{1/2} H_{3} + \cdots .
\end{equation}

To summarize, the LSWT Hamiltonian can be expressed in the following form:
\begin{align}
    H^{(2)} 
    = & \frac{1}{2}\sum_{\mathbf{k}} \Big(     
        \Psi_{\mathbf{k}}^\dagger \hat{\mathsf{H}}_{\mathbf{k}}^{\phantom{\dagger}} \Psi_{\mathbf{k}}^{\phantom{\dagger}} 
    - \Tr [\mathsf{A}_{\mathbf{k}}] \Big)
\end{align}
where $\Psi^\dagger_{\mathbf{k}}=\left( a^\dagger_{\mathbf{k}}, b^\dagger_{\mathbf{k}},\dots, a^{\phantom \dagger}_{-\mathbf{k}},b^{\phantom \dagger}_{-\mathbf{k}},\dots\right)$ is a vector of length $2 M_s$, with $M_s$ being the number of magnetic sublattices, and $\hat{\mathsf{H}}_{\mathbf{k}}$ is a $2M_s\times 2M_s$ matrix given by:
\begin{equation}
    \hat{\mathsf{H}}_{\mathbf{k}} = 
    \begin{pmatrix}
        \hat{\mathsf{A}}^{\phantom \dagger}_{\mathbf{k}} &  \hat{\mathsf{B}}^{\phantom \dagger}_{\mathbf{k}}\\[0.5ex] 
        \hat{\mathsf{B}}^\dagger_{\mathbf{k}}  & \hat{\mathsf{A}}^*_{-\mathbf{k}}
    \end{pmatrix},
\end{equation}
where $\hat{\mathsf{A}}^{\dagger}_{\mathbf{k}} = \hat{\mathsf{A}}^{\phantom \dagger}_{\mathbf{k}}$ are Hermitian, which follows from the Hermiticity of the Hamiltonian, and $\hat{\mathsf{B}}^{\phantom *}_{-\mathbf{k}} = \hat{\mathsf{B}}^{\rm T}_{\mathbf{k}}$, which follows from the bosonic commutation relations. 

Finally, we provide explicit expressions for $\mathsf{H}_{\mathbf{k}}$ for a two-sublattice system $(\rho, \sigma) = (a, b)$ for the reader's convenience.
Let us first consider the case where $a \neq b$, which can be written in matrix form as: 

\textcolor{blue}{(SJ: I think $t^{--}_{ij}$ should be replaced by $(t^{--}_{ij})^{*}$ in equation 46, to satisfy the constraint on $B$, $\hat{\mathsf{B}}^{\phantom *}_{-\mathbf{k}} = \hat{\mathsf{B}}^{\rm T}_{\mathbf{k}} $) }

\textcolor{red}{ (SP): 
    I agree and this is a typo. The code is implemented without this typo.
}
\begin{equation}
    \mathsf{A}_{\mathbf{k}} = 
    \begin{pmatrix}
        \widetilde{h}_{a} - \sum_{\mathsf{L}(i,j)} \widetilde{J}_{ij}^{00} S_{b} \quad
        & \sum_{\mathsf{L}(i,j)}t^{+-}_{ij} e^{-i \mathbf{k} \cdot \bm{\delta}_{ij}} \\[0.5em]
        \sum_{\mathsf{L}(i,j)}t^{-+}_{ij} e^{i \mathbf{k} \cdot \bm{\delta}_{ij}} \quad &
        \widetilde{h}_{b} - \sum_{\mathsf{L}(i,j)} \widetilde{J}_{ij}^{00} S_{a}
    \end{pmatrix},
    \quad  \mathsf{B}_{\mathbf{k}} =
    \begin{pmatrix}
        0   & \sum_{\mathsf{L}(i,j)}t^{++}_{ij} e^{-i \mathbf{k} \cdot \bm{\delta}_{ij}} \\[0.5em]
        \sum_{\mathsf{L}(i,j)}t^{++}_{ij} e^{i \mathbf{k} \cdot \bm{\delta}_{ij}} &    0
    \end{pmatrix},
\end{equation}
If the two sublattices are identical ($a = b$), the matrices in the Hamiltonian are given as: \textcolor{blue}{(SJ: I think the factor 2 in $A_k$ (in equation 47) is unnecessary; the chemical potential is doublely counted for $A_k$ and $A_{-k}$, and divided by 2 in total LSWT Hamiltonian )}
\begin{equation}
    \begin{split}
        \mathsf{A}_{\mathbf{k}} 
        = & 2 \Big( \widetilde{h}_{a} - \sum_{\mathsf{L}(ij)} \widetilde{J}_{ij}^{00} S_{a} \Big)
        + \sum_{\mathsf{L}(i,j)}(t_{ij}^{+-} e^{-i \mathbf{k} \cdot \bm{\delta}_{ij}} + t^{-+}_{ij} e^{+i \mathbf{k} \cdot \bm{\delta}_{ij}}), \\
        \mathsf{B}_{\mathbf{k}} 
        = & \sum_{\mathsf{L}(i,j)}t^{++}_{ij} (e^{-i \mathbf{k} \cdot \bm{\delta}_{ij}} + e^{+i \mathbf{k} \cdot \bm{\delta}_{ij}})
        = 2 \sum_{\mathsf{L}(i,j)}t^{++}_{ij} \cos ( \mathbf{k} \cdot \bm{\delta}_{ij}) 
    \end{split}
\end{equation}
This momentum-space formulation simplifies the analysis of bosonic excitations and provides a basis for studying dispersion relations and system stability.



\subsection{Diagonalization of Quadratic Boson Hamiltonian}

In this section, we discuss the diagonalization of the boson Hamiltonian. 
The quadratic boson Hamiltonian can be diagonalized by the Bogoliubov transformation. 
Consider the following Bogoliubov transformation:
\begin{equation}
    \hat{b}_{\mathbf{k}, \mu} 
    = \big[ \mathsf{P}_{\mathbf{k}}^{\phantom \dagger} \hat{\bm \beta}_{\mathbf{k}}^{\phantom \dagger} + \mathsf{Q}_{-\mathbf{k}}^{\phantom \dagger} \hat{\bm \beta}_{-\mathbf{k}}^{\dagger} \big]_{\mu}
    = \sum_{\nu = 1}^{M_s} [\mathsf{P}_{\mathbf{k}}]_{\mu\nu}^{\phantom \dagger} \hat{\beta}_{\mathbf{k}, \nu}^{\phantom \dagger} + [\mathsf{Q}_{-\mathbf{k}}]_{\mu\nu}^{\phantom \dagger} \hat{\beta}_{-\mathbf{k}, \nu}^{\dagger},
\end{equation}
where the indices $\mu, \nu$ are defined over $\mu, \nu = 1, \cdots, M_s$.
In order to preserve the bosonic commutation relations, we require the following identities:
\begin{align*}
    [\hat{b}_{  \mathbf{k}, \mu}^{\phantom \dagger}, \hat{b}_{  \mathbf{q}, \nu}^{\dagger}] 
    = & \big[ \mathsf{P}_{\mathbf{k}}\mathsf{P}_{\mathbf{k}}^{\dagger} - \mathsf{Q}_{-\mathbf{k}}\mathsf{Q}_{-\mathbf{k}}^{\dagger} \big]_{\mu\nu} \delta_{\mathbf{k},\mathbf{q}}
    = + \delta_{\mu,\nu}    \delta_{\mathbf{k}, \mathbf{q}} ,
    \qquad \quad
    [\hat{b}_{  \mathbf{k}, \mu}, \hat{b}_{-\mathbf{q}, \nu}] 
    = \big[ \mathsf{P}_{\mathbf{k}}\mathsf{Q}_{\mathbf{k}}^{\mathrm{T}} - \mathsf{Q}_{-\mathbf{k}}\mathsf{P}_{-\mathbf{k}}^{\mathrm{T}} \big]_{\mu\nu} \delta_{\mathbf{k},\mathbf{q}} = 0, \\
    [\hat{b}_{-\mathbf{k}, \mu}^{\dagger}, \hat{b}_{-\mathbf{q}, \nu}^{\phantom \dagger}] 
    = & \big[ \mathsf{P}_{-\mathbf{k}}^{*}\mathsf{P}_{-\mathbf{k}}^{\mathrm{T}} - \mathsf{Q}_{\mathbf{k}}^{*}\mathsf{Q}_{\mathbf{k}}^{\mathrm{T}} \big]_{\mu\nu} \delta_{\mathbf{k},\mathbf{q}}
    = - \delta_{\mu,\nu}    \delta_{\mathbf{k}, \mathbf{q}}, 
    \qquad \quad
    [\hat{b}_{-\mathbf{k}, \mu}^{\dagger}, \hat{b}_{  \mathbf{q}, \nu}^{\dagger}] 
    = \big[ \mathsf{Q}_{\mathbf{k}}^{*}\mathsf{P}_{\mathbf{k}}^{\dagger} - \mathsf{P}_{-\mathbf{k}}^{*}\mathsf{Q}_{-\mathbf{k}}^{\dagger} \big]_{\mu\nu} \delta_{\mathbf{k},\mathbf{q}}
    = 0.
\end{align*}
These identities can be represented in a compact matrix form:
\begin{equation}
    \begin{pmatrix}
       \mathsf{P}_{\mathbf{k}} & \mathsf{Q}_{-\mathbf{k}} \\ 
       \mathsf{Q}_{\mathbf{k}}^{*} & \mathsf{P}_{-\mathbf{k}}^{*}
    \end{pmatrix}
    \begin{pmatrix}
        1 & 0 \\ 0 & -1
    \end{pmatrix}
    \begin{pmatrix}
       \mathsf{P}_{\mathbf{k}}^{\dagger} & \mathsf{Q}_{\mathbf{k}}^{\rm T} \\ 
       \mathsf{Q}_{-\mathbf{k}}^{\dagger} & \mathsf{P}_{-\mathbf{k}}^{\rm T}
    \end{pmatrix} 
    = 
    \begin{pmatrix} 
        \mathsf{P}_{\mathbf{k}}\mathsf{P}_{\mathbf{k}}^{\dagger} - \mathsf{Q}_{-\mathbf{k}}\mathsf{Q}_{-\mathbf{k}}^{\dagger} & 
        \mathsf{P}_{\mathbf{k}}\mathsf{Q}_{\mathbf{k}}^{\rm T}   - \mathsf{Q}_{-\mathbf{k}}\mathsf{P}_{-\mathbf{k}}^{\rm T}     \\
        \mathsf{Q}_{\mathbf{k}}^{*}\mathsf{P}_{\mathbf{k}}^{\dagger} - \mathsf{P}_{-\mathbf{k}}^{*}\mathsf{Q}_{-\mathbf{k}}^{\dagger} & 
            \mathsf{Q}_{\mathbf{k}}^{*}\mathsf{Q}_{\mathbf{k}}^{\rm T} - \mathsf{P}_{-\mathbf{k}}^{*}\mathsf{P}_{-\mathbf{k}}^{\rm T}
    \end{pmatrix}
    = 
    \begin{pmatrix}
        1 & 0 \\ 0 & -1
    \end{pmatrix}
\end{equation}

By adopting the above convention (i.e., the transformation of $\hat{\bm b}_{\mathbf{k}} \rightarrow \hat{\bm b}^{\dagger}_{-\mathbf{k}}$), the (para)-unitary transformation matrix $\mathsf{T}_{\mathbf{k}}$ has the following form: 
\begin{equation}
    \Psi_{\mathbf{k}} = \mathsf{T}_{\mathbf{k}} \widetilde{\Psi}_{\mathbf{k}}: \quad
    \begin{pmatrix}
        \hat{\bm b}_{\mathbf{k}} \\
        \hat{\bm b}_{-\mathbf{k}}^{\dagger}
    \end{pmatrix}
    =
    \begin{pmatrix}
        \mathsf{P}_{\mathbf{k}}       & \mathsf{Q}_{-\mathbf{k}} \\ 
        \mathsf{Q}_{\mathbf{k}}^{*}   & \mathsf{P}_{-\mathbf{k}}^{*}
    \end{pmatrix}
    \begin{pmatrix}
        \hat{\bm \beta}_{ \mathbf{k}} \\
        \hat{\bm \beta}_{-\mathbf{k}}^{\dagger}
    \end{pmatrix}.
\end{equation}
Suppose the unitary transformation $\mathsf{T}_{\mathbf{k}}$ diagonalizes the Hamiltonian matrix $\mathsf{H}_{\mathbf{k}}$:
\begin{equation}
    \mathsf{T}_{\mathbf{k}}^{\dagger} \mathsf{H}_{\mathbf{k}}^{\phantom \dagger} \mathsf{T}_{\mathbf{k}} = \mathsf{E}_{\mathbf{k}}: \qquad
        \begin{pmatrix}
        \mathsf{P}_{\mathbf{k}}^{\dagger}     & \mathsf{Q}_{  \mathbf{k}}^{\rm T} \\[0.3em] 
        \mathsf{Q}_{-\mathbf{k}}^{\dagger}    & \mathsf{P}_{- \mathbf{k}}^{\rm T}
    \end{pmatrix}
    \begin{pmatrix}
        \mathsf{A}^{\phantom \dagger}_{\mathbf{k}}      &   \mathsf{B}^{\phantom \dagger}_{\mathbf{k}}\\[0.3em] 
        \mathsf{B}^\dagger_{\mathbf{k}}                 &   \mathsf{A}^*_{-\mathbf{k}}
    \end{pmatrix}
    \begin{pmatrix}
        \mathsf{P}_{\mathbf{k}}       & \mathsf{Q}_{-\mathbf{k}} \\[0.3em] 
        \mathsf{Q}_{\mathbf{k}}^{*}   & \mathsf{P}_{-\mathbf{k}}^{*}
    \end{pmatrix}
    =
    \begin{pmatrix}
        \mathsf{E}_{\mathbf{k}} & 0 \\ 0 & \widetilde{\mathsf{E}}_{-\mathbf{k}}
    \end{pmatrix},
\end{equation}
The two diagonal matrices $\mathsf{E}_{\mathbf{k}}$ and $\widetilde{\mathsf{E}}_{-\mathbf{k}}$ are real-valued, and their explicit forms are written as:
\begin{align*}
    \mathsf{E}_{\mathbf{k}} = & 
    \mathsf{P}_{\mathbf{k}}^{\dagger} \mathsf{A}_{\mathbf{k}} \mathsf{P}_{\mathbf{k}} + 
    \mathsf{P}_{\mathbf{k}}^{\dagger} \mathsf{B}_{\mathbf{k}} \mathsf{Q}_{\mathbf{k}}^{*} + 
    \mathsf{Q}_{\mathbf{k}}^{\rm T} \mathsf{B}_{\mathbf{k}}^{\dagger} \mathsf{P}_{\mathbf{k}} + 
    \mathsf{Q}_{\mathbf{k}}^{\rm T} \mathsf{A}_{-\mathbf{k}}^{*} \mathsf{Q}_{\mathbf{k}}^{*},
    \\
    \widetilde{\mathsf{E}}_{-\mathbf{k}} = & 
    \mathsf{Q}_{-\mathbf{k}}^{\dagger} \mathsf{A}_{\mathbf{k}} \mathsf{Q}_{-\mathbf{k}} + 
    \mathsf{Q}_{-\mathbf{k}}^{\dagger} \mathsf{B}_{\mathbf{k}} \mathsf{P}_{-\mathbf{k}}^{*} + 
    \mathsf{P}_{-\mathbf{k}}^{\rm T} \mathsf{B}_{\mathbf{k}}^{\dagger} \mathsf{Q}_{-\mathbf{k}} + 
    \mathsf{P}_{-\mathbf{k}}^{\rm T} \mathsf{A}_{-\mathbf{k}}^{*} \mathsf{P}_{-\mathbf{k}}^{*}.
\end{align*}
Note that the two diagonal matrices can be mapped to each other under conjugation and momentum inversion. 
Thus, we obtain the energy bands ($E_{\mathbf{k}}$ and $E_{-\mathbf{k}}$) when we diagonalize the quadratic boson Hamiltonian via Colpa's method~\cite{Colpa1978Diagonalization}:
\begin{equation}
    \widetilde{\mathsf{E}}_{-\mathbf{k}} 
    \xrightarrow[\mathbf{k} \rightarrow -\mathbf{k}]{\rm conj} 
    \widetilde{\mathsf{E}}_{\mathbf{k}}^{*}:  \qquad
    \widetilde{\mathbf{E}}_{\mathbf{k}}^{*} =
    \mathsf{P}_{\mathbf{k}}^{\dagger}   \mathsf{A}_{\mathbf{k}}             \mathsf{P}_{\mathbf{k}}   +
    \mathsf{P}_{\mathbf{k}}^{\dagger}   \mathsf{B}^{\rm T}_{-\mathbf{k}}    \mathsf{Q}_{\mathbf{k}}^{*} +
    \mathsf{Q}_{\mathbf{k}}^{\rm T}     \mathsf{B}^{*}_{-\mathbf{k}}        \mathsf{P}_{\mathbf{k}}     +
    \mathsf{Q}_{\mathbf{k}}^{\rm T}     \mathsf{A}_{-\mathbf{k}}^{*}        \mathsf{Q}_{\mathbf{k}}^{*}
    = \mathbf{E}_{\mathbf{k}},
\end{equation}
where $\mathsf{B}^{\rm T}_{-\mathbf{k}} = \mathsf{B}_{\mathbf{k}}^{\phantom T}$ and $\mathsf{B}^{*}_{-\mathbf{k}} = \mathsf{B}_{\mathbf{k}}^{\dagger}$.
Since $\mathsf{E}_{\mathbf{k}}$ is a \emph{real} diagonal matrix, we conclude $\widetilde{\mathsf{E}}_{\mathbf{k}} = \mathsf{E}_{\mathbf{k}}$. 
Thus, the canonical form of the quadratic bosonic Hamiltonian reads:
\begin{align}
    \hat{H}_{2}
    = & \frac{1}{2} \sum_{\mathbf{k}} 
    \Big( (\widetilde{\Psi}^{\dagger}_{\mathbf{k}} \mathsf{T}_{\mathbf{k}}^{\dagger}) \mathsf{H}_{\mathbf{k}}^{\phantom \dagger} (\mathsf{T}_{\mathbf{k}}^{\phantom \dagger} \widetilde{\Psi}_{\mathbf{k}}^{\phantom \dagger} )- \Tr [\mathsf{A}_{\mathbf{k}}] \Big) 
    = \frac{1}{2} \sum_{\mathbf{k}} 
    \Big( \widetilde{\Psi}^{\dagger}_{\mathbf{k}} \mathsf{E}_{\mathbf{k}} \widetilde{\Psi}_{\mathbf{k}} - \Tr [\mathsf{A}_{\mathbf{k}}] \Big) \nonumber \\
    = & \frac{1}{2} \sum_{\mathbf{k}} 
    \begin{pmatrix}
        \hat{\bm \beta}_{\mathbf{k}} \\
        \hat{\bm \beta}_{-\mathbf{k}}^{\dagger}
    \end{pmatrix}^{\dagger}
    \begin{pmatrix}
        \mathbf{E}_{\mathbf{k}} & 0 \\ 0 & \widetilde{\mathbf{E}}_{-\mathbf{k}}
    \end{pmatrix}
    \begin{pmatrix}
        \hat{\bm \beta}_{\mathbf{k}} \\
        \hat{\bm \beta}_{-\mathbf{k}}^{\dagger}
    \end{pmatrix} 
    - \frac{1}{2} \Tr [\mathsf{A}_{\mathbf{k}}]
    \nonumber \\
    = & \frac{1}{2} \sum_{\mathbf{k}, \sigma} 
    \big( 
        E_{\mathbf{k}, \sigma} \hat{\beta}^{\dagger}_{\mathbf{k}, \sigma}   \hat{\beta}^{\phantom \dagger}_{\mathbf{k}, \sigma} +
        E_{-\mathbf{k}, \sigma} \hat{\beta}^{\phantom \dagger}_{-\mathbf{k}, \sigma} \hat{\beta}^{\dagger}_{-\mathbf{k}, \sigma}
    \big) 
    - \frac{1}{2} \Tr [\mathsf{A}_{\mathbf{k}}] \\
    = & \sum_{\mathbf{k}, \sigma} E_{\mathbf{k}, \sigma} \hat{\beta}_{\mathbf{k}, \sigma}^{\dagger} \hat{\beta}_{\mathbf{k}, \sigma}^{\phantom \dagger} 
    + \frac{1}{2} \sum_{\mathbf{k}, \sigma} \big( E_{\mathbf{k}, \sigma} - \Tr [\mathsf{A}_{\mathbf{k}}] \big),
\end{align}
where $E_{\mathbf{k}} \ge 0$ for all $\mathbf{k} \in \mathrm{BZ}$. 
The constant term is called the ``zero-point energy,'' as it represents the ground state energy of the quantum system at zero temperature (the Gibbs state at $T=0$).
This gives the quantum correction to the classical energy in the LSWT framework:
\begin{equation}\label{eq: def. of zero point energy}
    E_{0} = E_{\rm cl} + \frac{1}{2} \sum_{\mathbf{k} \in \mathrm{FBZ}} \Big( \sum_{\sigma=a,b,\cdots} E_{\mathbf{k}, \sigma} -  \Tr[\mathsf{A}_{\mathbf{k}}] \Big),
\end{equation}
and the energy density per spin site is:
\begin{equation}
    \mathcal{E} 
    = \frac{E_{0}}{N} 
    = \frac{1}{N} \Big[ E_{\rm cl} + \frac{1}{2} \sum_{\mathbf{k} \in \mathrm{FBZ}} \Big( \sum_{\sigma=a,b,\cdots} E_{\mathbf{k}, \sigma} -  \Tr[\mathsf{A}_{\mathbf{k}}] \Big) \Big],
\end{equation}
where $N=Lm_{s}$ is the total number of spin sites in the system, with $L$ being the number of unit cells and $m_s$ the number of magnetic sublattices per unit cell.



\subsection{Paraunitary Diagonalization}
Now, the remaining problem is to find a paraunitary matrix $\mathsf{T}_{\kvec}$ (we will suppress the $\kvec$ subscript in this section for clarity, so $\mathsf{H}_{\kvec} \rightarrow \mathsf{H}$, etc.) such that
\begin{equation}
    \mathsf{T}^{\dagger}\mathsf{HT} = \mathsf{E}, \quad \text{with}
    \quad \mathsf{TJT}^{\dagger} = \mathsf{J},
\end{equation}
where $\mathsf{E}$ is a diagonal matrix containing the magnon energies, and $\mathsf{J}$ is the symplectic metric tensor. The paraunitary transformation $\mathsf{T}$ exists if and only if $\mathsf{H}$ is positive definite.

The explicit form of the paraunitary matrix is given by 
\begin{equation}
    \mathsf{T} = (\mathsf{K}^{\dagger})^{-1} \mathsf{V} \mathsf{(\Lambda J)}^{1/2},
\end{equation}
where $\mathsf{K}$ comes from the Cholesky decomposition of the Hamiltonian. Since the Hamiltonian matrix $\mathsf{H}$ is positive definite, the Cholesky decomposition exists, satisfying:
\begin{equation}
    \mathsf{H} = \mathsf{KK}^{\dagger},
\end{equation}
and $\mathsf{V}$ is the unitary matrix that diagonalizes $\mathsf{K}^{\dagger} \mathsf{JK}$:
\begin{equation}
    \mathsf{V \Lambda V}^{\dagger} = \mathsf{K}^{\dagger} \mathsf{JK},
\end{equation}
where $\mathsf{\Lambda}$ is a diagonal matrix.

We will now verify that this construction satisfies the required paraunitary conditions. First, let us check $\mathsf{TJT}^{\dagger} = \mathsf{J}$:
\begin{align*}
    \mathsf{TJT}^{\dagger} &= 
    \big( (\mathsf{K}^{\dagger})^{-1} \mathsf{V} \mathsf{(\Lambda J)}^{1/2} \big) \mathsf{J} \big( \mathsf{(\Lambda J)}^{1/2} \big)^{\dagger} \mathsf{V}^{\dagger} \mathsf{K}^{-1} \\
    &= (\mathsf{K}^{\dagger})^{-1} \mathsf{V} \mathsf{(\Lambda J)}^{1/2} \mathsf{J} \mathsf{(J\Lambda)}^{1/2} \mathsf{V}^{\dagger} \mathsf{K}^{-1} \\
    &= (\mathsf{K}^{\dagger})^{-1} \mathsf{V} \mathsf{\Lambda} \mathsf{V}^{\dagger} \mathsf{K}^{-1} \\
    &= (\mathsf{K}^{\dagger})^{-1} (\mathsf{K}^{\dagger} \mathsf{JK}) \mathsf{K}^{-1} = \mathsf{J}
\end{align*}
where we can also verify that $\mathsf{T}^{\dagger} \mathsf{JT} = \mathsf{J}$.
% \begin{align*}
%     \mathsf{T}^{\dagger} \mathsf{JT} &= 
%     \big( \mathsf{(\Lambda J)}^{1/2} \big)^{\dagger} \mathsf{V}^{\dagger} \mathsf{K}^{-1} \mathsf{J} (\mathsf{K}^{\dagger})^{-1} \mathsf{V} \mathsf{(\Lambda J)}^{1/2} \\
%     &= \mathsf{(J\Lambda)}^{1/2} \mathsf{V}^{\dagger} \mathsf{K}^{-1} \mathsf{J} (\mathsf{K}^{\dagger})^{-1} \mathsf{V} \mathsf{(\Lambda J)}^{1/2} \\
%     &= \mathsf{(J\Lambda)}^{1/2} \mathsf{V}^{\dagger} (\mathsf{K}^{\dagger} \mathsf{JK})^{-1} \mathsf{V} \mathsf{(\Lambda J)}^{1/2} \\
%     &= \mathsf{(J\Lambda)}^{1/2} \mathsf{V}^{\dagger} (\mathsf{V \Lambda V}^{\dagger})^{-1} \mathsf{V} \mathsf{(\Lambda J)}^{1/2} \\
%     &= \mathsf{(J\Lambda)}^{1/2} \mathsf{V}^{\dagger} \mathsf{V} \mathsf{\Lambda}^{-1} \mathsf{V}^{\dagger} \mathsf{V} \mathsf{(\Lambda J)}^{1/2} \\
%     &= \mathsf{(J\Lambda)}^{1/2} \mathsf{\Lambda}^{-1} \mathsf{(\Lambda J)}^{1/2} \\
%     &= \mathsf{J}
% \end{align*}

Next, let us check that $\mathsf{T}$ diagonalizes the Hamiltonian as required:
\begin{align*}
    \mathsf{T}^{\dagger} \mathsf{HT}
    &= \big( \mathsf{(\Lambda J)}^{1/2} \big)^{\dagger} \mathsf{V}^{\dagger} \mathsf{K}^{-1} \cdot \mathsf{KK}^{\dagger} \cdot (\mathsf{K}^{\dagger})^{-1} \mathsf{V} \mathsf{(\Lambda J)}^{1/2} \\
    &= \mathsf{(J\Lambda)}^{1/2} \mathsf{V}^{\dagger} \mathsf{V} \mathsf{(\Lambda J)}^{1/2} \\
    &= \mathsf{(J\Lambda)}^{1/2} \mathsf{(\Lambda J)}^{1/2} \\
    &= \mathsf{J\Lambda} = \mathsf{\Lambda J} = \mathsf{E}
\end{align*}
where we identify the diagonal energy matrix as $\mathsf{E} = \mathsf{J\Lambda} = \mathsf{\Lambda J}$. This equality holds because both $\mathsf{J}$ and $\mathsf{\Lambda}$ are diagonal matrices, with $\mathsf{J}$ having half positive and half negative entries, and $\mathsf{\Lambda}$ arranged to ensure that $\mathsf{E}$ is always positive definite.



\newpage



\section{Physical Quantities in Linear Spin Wave Theory}

In this Section, I summarize the methods for obtaining important physical quantities in the frame of linear spin wave theory.
The physical quantities computable in LSWT are the following:


\begin{table*}[!h]
    \centering
    \begin{tabular}{c|c|c|c}
    \hline\hline
    \textbf{Quantity} & \textbf{Definition} & \textbf{Calculation in LSWT} & \textbf{Sect.} \\
    \hline
    \multicolumn{3}{c|}{\bf Thermodynamics} & \textbf{III} \\[0.5ex]
    \hline
    $Z_{\beta}$ % partition function
    & 
    $\text{Tr} [e^{-\beta \hat{H}}]$ 
    & $e^{-\beta E_{0}} \prod_{\mathbf{k}, \mu} \frac{1}{1 - e^{-\beta E_{\mathbf{k}, \mu}}}$ 
    & \ref{subsec: Partition Function}
    \\[1ex]
    $U$ 
    & $\langle \hat{H} \rangle$ & $\sum_{\mathbf{k}, \mu} E_{\mathbf{k}, \mu} n_{\mathbf{k}, \mu}$
    & \ref{subsec: Internal Energy}
    \\[1ex]
    $F$ 
    & $- k_B T \ln Z_{\beta}$ 
    & $  k_B T \sum_{\mathbf{k}, \mu} \ln \left( 1 - e^{-\beta E_{\mathbf{k}, \mu}} \right) + E_{0} $ 
    & \ref{subsec: Free Energy}
    \\[1ex]
    $S$ 
    & $-k_{\rm B}\Tr \big( \rho \ln \rho \big)$
    & $k_B \sum_{\mathbf{k}, \mu} \left[ (1 + n_{\mathbf{k}, \mu}) \ln (1 + n_{\mathbf{k}, \mu}) - n_{\mathbf{k}, \mu} \ln n_{\mathbf{k}, \mu} \right]$ 
    & \ref{subsec: Entropy}
    \\[1ex]
    $C$ 
    & $\frac{\partial U}{\partial T}$ 
    & $\sum_{\mathbf{k}, \mu} \frac{\partial n_{\mathbf{k}, \mu}}{\partial T} E_{\mathbf{k}, \mu}$ 
    & \ref{subsec: Specific Heat} 
    \\[1ex]
    $\langle \hat{n}_\mu \rangle$ 
    & $\sum_{\mathbf{k}} \langle a^{\dagger}_{\mathbf{k}, \mu} a_{\mathbf{k}, \mu} \rangle$  
    & $\sum_{\mathbf{k}} [\langle \Psi_{\mathbf{k}}^{\phantom\dagger} \Psi_{\mathbf{k}}^{\dagger} \rangle]_{m_{s} + \mu, m_{s} + \mu}$ 
    & \ref{subsec: Number expectation value/Spin moment} \\
    [1ex]
    $\langle \widetilde{S}_\mu \rangle$ 
    & $S_{\mu} - \langle \hat{n}_\mu \rangle$  
    & $S_{\mu} - \frac{1}{L} \sum_{\mathbf{k}} \langle \hat{n}_{\mathbf{k}, \mu} \rangle$ & \ref{subsec: Number expectation value/Spin moment} \\
    [1ex]
    \hline
    \multicolumn{3}{c|}{\bf Correlations} & \textbf{IV} \\[0.5ex]
    \hline
    $\mathcal{C}^{\alpha \beta}_{\mu\nu}(\mathbf{k}, t)$ 
    & 
    $\langle S_{\mathbf{k},\mu}^\alpha(t) S_{-\mathbf{k},\nu}^\beta(0) \rangle$ 
    & $\sum_{(m, n)_{\mu\nu}}
    [\mathsf{U}^\alpha \mathsf{S}_{\mathbf{k}} \mathsf{T}_{\mathbf{k}} \mathsf{N}_{\mathbf{k}}(t) \mathsf{T}_{\mathbf{k}}^\dagger \bar{\mathsf{S}}_{\mathbf{k}} \bar{\mathsf{U}}^{\beta}]_{m n}$ 
    & \ref{subsec: Real-Time (Dynamical) Spin-Spin Correlation Function} \\[1ex]
    $\mathcal{S}^{\alpha \beta}(\mathbf{k},\omega)$ 
    & $\frac{1}{2 \pi} \int dt \, e^{i \omega t} \mathcal{C}^{\alpha \beta}(\mathbf{k}, t)$ 
    & $\sum_{(m, n)_{\mu\nu}} [ \mathsf{U}^\alpha \mathsf{S}_{\mathbf{k}} \mathsf{T}_{\mathbf{k}} \mathcal{F}[\mathsf{N}_{\mathbf{k}}](\omega; \eta) \mathsf{T}_{\mathbf{k}}^\dagger \bar{\mathsf{S}}_{\mathbf{k}} \bar{\mathsf{U}}^{\beta}]_{m, n}$ 
    & \ref{subsec: Structure Factor} \\[1ex]
    $\mathcal{A}^{\alpha \beta}(\mathbf{k}, \omega)$ 
    & $-\frac{1}{\pi} \operatorname{Im} [G_{\mathsf{R}}^{\alpha \beta}(\mathbf{k}, \omega)]$ & $-\frac{2}{\pi} \operatorname{Im} \sum_{m, n} [\mathsf{U}^\alpha \mathsf{S}_{\mathbf{k}} \mathsf{T}_{\mathbf{k}} \mathcal{G}[\mathsf{N}_{\mathbf{k}}] (\omega; \eta)
    \mathsf{T}_{\mathbf{k}}^\dagger \mathsf{S}_{\mathbf{k}} \mathsf{U}^\beta]_{m, n}$ & \ref{subsec: Spectral Function} \\
    \hline
    \multicolumn{3}{c|}{\bf Topology} & \textbf{V} \\[0.5ex]
    \hline
    $Q_{\rm skyr}$ 
    & $\frac{1}{4\pi} \sum_{\triangle} \chi_{\triangle}$ 
    & $\chi_{\triangle} = 2 \arctan \left( \frac{|\mathbf{m}_i \cdot (\mathbf{m}_j \times \mathbf{m}_k)|}{1 + \mathbf{m}_i \cdot \mathbf{m}_j + \mathbf{m}_j \cdot \mathbf{m}_k + \mathbf{m}_k \cdot \mathbf{m}_i} \right)$
    & \ref{subsec: Skyrmion number}
    \\
    $C_n$ & $\frac{1}{2 \pi} \int_{\mathrm{FBZ}} \Omega_{n, \mathbf{k}} \, d^2 \mathbf{k}$ 
    & $\Omega_{n, \mathbf{k}} = -2 \operatorname{Im} \sum_{m \neq n} \frac{[\mathsf{J} \mathsf{T}_{\mathbf{k}}^\dagger (\partial_x \mathsf{H}_{\mathbf{k}}) \mathsf{T}_{\mathbf{k}}]_{nm} [\mathsf{J} \mathsf{T}_{\mathbf{k}}^\dagger (\partial_y \mathsf{H}_{\mathbf{k}}) \mathsf{T}_{\mathbf{k}}]_{mn}}{(E_{\mathbf{k}, nn} - E_{\mathbf{k}, mm})^2}$ 
    & \ref{subsec: Chern number} \\[1ex]
    $\kappa_{xy}$ 
    & $\text{Thermal Hall conductivity}$
    & $ - \frac{k_{B}^2 T}{V} \sum_{n} \int_{\mathbf{k}} d^{2} \mathbf{k} \left( c_2(n_{\mathbf{k}, n}) - \frac{\pi^2}{3} \right) \Omega_{n, \mathbf{k}}$ 
    &  \ref{subsec: Thermal Hall Conductance} \\
    \hline\hline
    \end{tabular}
    \caption{Physical Quantities in Linear Spin Wave Theory (LSWT). 
    Here, $\beta = 1/k_{B}T$,
    $n_{\mathbf{k}, \mu} = \frac{1}{e^{\beta E_{\mathbf{k}, \mu}} - 1}$ is the Bose-Einstein distribution function,
    $\Omega_{n, \mathbf{k}}$ is the Berry curvature, and 
    $c_2(x)$ is the Spence function.
    }
    \label{tab:lswt_quantities}
\end{table*}









\newpage
\section{Thermodynamics in Linear Spin Wave Theory}

\subsection{Partition function}\label{subsec: Partition Function}
In this section, we derive the partition function for linear spin wave theory. 
Recall that the Hamiltonian for the spin system with normal bosons is written as 
\begin{align}
    \hat{H} = E_{\rm cl} + \sum_{\kvec, \mu} E_{\kvec, \mu} \Big( \beta^{\dagger}_{\kvec, \mu} \beta_{\kvec, \mu} + \frac{1}{2} \Big) - \frac{1}{2} \Tr [\mathsf{H}_{\kvec}] 
    = E_{0} + \sum_{\kvec \in \mathrm{FBZ}} \sum_{\mu = a, b, \cdots} E_{\kvec, \mu} \hat{n}_{\kvec, \mu},
\end{align}
where $E_{0} = E_{\rm cl} + \frac{1}{2} \sum_{\kvec} \big( \sum_{\mu} E_{\kvec, \mu} - \Tr [\mathsf{A}_{\kvec}] \big)$ represents the zero-point energy [Eqn.~\eqref{eq: def. of zero point energy}] and $\hat{n}_{\mu}(\kvec) = \beta^{\dagger}_{\kvec, \mu} \beta_{\kvec, \mu}$ is the magnon number operator. Here, $E_{\kvec, \mu}$ are the magnon energy eigenvalues with band index $\mu$.

The partition function is given by
\begin{equation}
    Z_{\beta} \equiv \Tr \big( e^{- \beta \hat{H}} \big) 
    = e^{- \beta E_{0}}  \Tr \big( e^{- \beta \sum_{\kvec, \mu} E_{\kvec, \mu} \beta^{\dagger}_{\kvec, \mu} \beta_{\kvec, \mu}} \big) 
    = e^{- \beta E_{0}} \Tr \big( e^{- \beta \sum_{\kvec, \mu} E_{\kvec, \mu} \hat{n}_{\kvec, \mu}} \big),
\end{equation}
Here, the constant term $E_0$ contributes to the internal energy but is not important for thermodynamic quantities that involve derivatives of $\ln \tilde{Z}_{\beta}$, where $\tilde{Z}_{\beta}$ represents the partition function without the zero-point energy contribution.

The partition function $Z_{\beta}$ can be evaluated as
\begin{subequations}\label{eq: explicit expressions for Z and log-Z}
    \begin{align}
        \Tr \big( e^{- \beta \sum_{\kvec, \mu} E_{\kvec, \mu} \hat{n}_{\mu}(\kvec)} \big) 
        = & \Tr \Big[ \prod_{\kvec, \mu} e^{- \beta E_{\kvec, \mu} \hat{n}_{\kvec, \mu}} \Big]
        = \prod_{\kvec, \mu} \Tr \Big[ e^{- \beta E_{\kvec, \mu} \hat{n}_{\kvec, \mu}} \Big] 
        =  \prod_{\kvec, \mu}  \frac{1}{1 - e^{- \beta E_{\kvec, \mu}}} \\
        - \ln Z_{\beta} &= \beta E_{0} + \sum_{\kvec, \mu} \ln \big( 1 - e^{- \beta E_{\kvec, \mu}} \big).
    \end{align}
\end{subequations}
The final expression for $-\ln Z_{\beta}$ is particularly useful for calculating thermodynamic quantities such as free energy, internal energy, specific heat, and entropy in the linear spin-wave approximation.


\subsection{Internal Energy}\label{subsec: Internal Energy}
The internal energy $U$ is the expectation value of the Hamiltonian $\hat{H}$, computed as:
\begin{equation}
    U \equiv \langle \hat{H} \rangle  = -\frac{\partial \ln Z_{\beta}}{\partial \beta},
\end{equation}
where $Z_{\beta} = e^{- \beta E_{0}} \tilde{Z}_{\beta}$, with $E_0$ being the zero-point energy and $\tilde{Z}_{\beta}$ the partition function excluding the zero-point energy contribution.

Computing the derivative, we get:
\begin{equation}
     - \frac{\partial \ln Z_{\beta}}{\partial \beta} 
    = E_{0} + \frac{\partial}{\partial \beta} \Big[  \sum_{\kvec, \mu} \ln \big( 1 - e^{- \beta E_{\kvec, \mu}} \big) \Big]
    = E_{0} + \sum_{\kvec, \mu} \frac{E_{\kvec, \mu} e^{-\beta E_{\kvec, \mu}}}{1 - e^{-\beta E_{\kvec, \mu}}}  
    = E_{0} + \sum_{\kvec, \mu} E_{\kvec, \mu} \frac{1}{e^{\beta E_{\kvec, \mu}} - 1},
\end{equation}
where we identify the Bose-Einstein distribution function:
\begin{equation}
    n_{\mu} (\kvec) = \frac{1}{e^{\beta E_{\kvec, \mu}} - 1},
\end{equation}
which represents the thermal expectation value of the magnon number operator $\langle\hat{n}_{\kvec, \mu}\rangle = \langle \hat{\beta}^{\dagger}_{\kvec, \mu} \hat{\beta}_{\kvec,\mu} \rangle$.
Therefore, the internal energy can be expressed as:
\begin{align}
    U &= E_0 + \sum_{\kvec, \mu} E_{\kvec, \mu} n_{\mu} (\kvec) \nonumber \\
    &= E_{\rm cl} + \frac{1}{2} \sum_{\kvec} \big( \sum_{\mu} E_{\kvec, \mu} - \Tr [\mathsf{A}_{\kvec}] \big) + \sum_{\kvec, \mu} E_{\kvec, \mu} n_{\mu}(\kvec)
\end{align}
This result expresses the internal energy as the sum of the classical ground state energy, the zero-point quantum fluctuations, and the thermal contribution from magnon excitations.



\subsection{Free Energy}\label{subsec: Free Energy}
The Helmholtz free energy ($F$) is defined as 
\begin{equation}\label{eq: def. of Helmholtz free energy}
    Z_{\beta} \equiv e^{-\beta F}
\end{equation}
where $Z_{\beta} = e^{- \beta E_0} \tilde{Z}_{\beta}$ is the partition function, with $\tilde{Z}_{\beta}$ being the partition function excluding the zero-point energy contribution.
From \eqref{eq: explicit expressions for Z and log-Z} and Eq.~\eqref{eq: def. of Helmholtz free energy}, we get:
\begin{align}
    F &= -\frac{1}{\beta} \ln(e^{- \beta E_0} \tilde{Z}_{\beta}) \nonumber \\
    &= E_0 - \frac{1}{\beta} \ln \tilde{Z}_{\beta} \nonumber \\
    &= E_{\rm cl} + \frac{1}{2} \sum_{\kvec} \big( \sum_{\mu} E_{\kvec, \mu} - \Tr [\mathsf{A}_{\kvec}] \big) + \frac{1}{\beta} \sum_{\kvec, \mu} \ln \left( 1 - e^{-\beta E_{\kvec, \mu}} \right),
\end{align}
where we use $-\ln \tilde{Z}_{\beta} = \sum_{\kvec, \mu} \ln \left( 1 - e^{-\beta E_{\kvec, \mu}} \right)$ from the partition function section, and $E_0 = E_{\rm cl} + \frac{1}{2} \sum_{\kvec} \big( \sum_{\mu} E_{\kvec, \mu} - \Tr [\mathsf{A}_{\kvec}] \big)$ is the zero-point energy.


\subsection{Entropy Expression}\label{subsec: Entropy}
In this section, we derive the expression for the thermal entropy in LSWT. 
The entropy is defined as 
\begin{equation}
    S \equiv -k_{\rm B}\Tr \big( \rho \ln \rho \big),
\end{equation}
where $\rho$ is the density operator, which is Gibbs states ($\rho = e^{-\beta \hat{H}}/Z_{\beta}$) especially in thermal equilibrium.
This definition coincides with the conventional definition of entropy in thermal physics:
\begin{align}
    - \frac{\partial F}{\partial T}
    = \frac{1}{T} \Big(-\frac{\partial}{\partial \beta} \ln Z_{\beta} \Big) - \big(-k_{\rm B} \ln Z_{\beta} \big)
    = \frac{U - F}{T}.
\end{align}
We can derive the same expression using Gibbs states:
\begin{align}
    S = & -k_{\rm B}\Tr \big( \rho \ln \rho \big)
    = -k_{\rm B}\Tr \big( \rho (-\beta \hat{H} - \ln Z_{\beta}) \big) \nonumber \\
    = &  \frac{ \Tr (\rho \hat{H}) - (-k_{\rm B}T\ln Z_{\beta})}{T} 
    = \frac{U - F}{T}.
\end{align}
Using our previous expressions for $U$ and $F$, we have:
\begin{align}
    U - F &= E_0 + \sum_{\kvec, \mu} E_{\kvec, \mu} n_{\mu}(\kvec) - \left( E_0 + \frac{1}{\beta} \sum_{\kvec, \mu} \ln \left( 1 - e^{-\beta E_{\kvec, \mu}} \right) \right) \nonumber \\ 
    &= \sum_{\kvec, \mu} E_{\kvec, \mu} n_{\mu}(\kvec) - \frac{1}{\beta} \sum_{\kvec, \mu} \ln \left( 1 - e^{-\beta E_{\kvec, \mu}} \right).
\end{align}

Since $S = k_{\rm B} \beta (U - F)$ (noting $\beta = \frac{1}{k_B T}$ and adjusting units appropriately):
\begin{equation}
    S = k_{\rm B} \sum_{\kvec, \mu} \beta E_{\kvec, \mu} n_{\mu} (\kvec) - k_{\rm B} \sum_{\kvec, \mu} \ln \left( 1 - e^{-\beta E_{\kvec, \mu}} \right).
\end{equation}

Now, we can rewrite this expression using the properties of the Bose-Einstein distribution function $n_{\kvec, \mu}$. From the definition $n_{\kvec, \mu} = \frac{1}{e^{\beta E_{\kvec, \mu}} - 1}$, we can derive:
\begin{equation}
    \beta E_{\kvec, \mu} = \ln \Big( \frac{1 + n_{\mu}(\kvec)}{n_{\mu}(\kvec)} \Big), 
    \quad \text{and} \quad
    \ln \left( 1 - e^{-\beta E_{\kvec, \mu}} \right) 
    = -\ln \big( 1 + n_{\mu} (\kvec) \big), 
\end{equation}
Substituting these relations into our entropy expression, we obtain:
\begin{equation}
    S = k_{\rm B} \sum_{\kvec, \mu} \big[ (1 + n_{\mu}(\kvec)) \ln (1 + n_{\mu}(\kvec)) - n_{\mu}(\kvec) \ln n_{\mu} (\kvec) \big]
\end{equation}
This is the standard entropy formula for a system of non-interacting bosons, which emerges naturally from our linear spin wave theory.




\subsection{Specific Heat}\label{subsec: Specific Heat}
The specific heat $C$ is defined as the temperature derivative of the internal energy:
\begin{equation}
    C = \frac{\partial U}{\partial T} = \sum_{\kvec, \mu} E_{\kvec, \mu} \frac{\partial n_{\mu}(\kvec)}{\partial T},
\end{equation}
where $n_{\mu}(\kvec)$ is the Bose-Einstein distribution for magnons $\mu$. 
The derivative of the distribution with respect to temperature is 
\begin{equation}
    \frac{\partial n_{\mu}(\kvec)}{\partial T} 
    =  \frac{\partial n_{\mu}(\kvec)}{\partial \beta} \frac{\partial \beta}{\partial T} \\
    =  \Big( - \frac{E_{\kvec, \mu} e^{\beta E_{\kvec, \mu}}}{\left( e^{\beta E_{\kvec, \mu}} - 1 \right)^2} \Big) \Big( - \frac{\beta}{T}\Big) 
    =  \frac{\beta}{T} \frac{E_{\kvec, \mu} e^{\beta E_{\kvec, \mu}}}{\left( e^{\beta E_{\kvec, \mu}} - 1 \right)^2}.
\end{equation}
Therefore, the specific heat is given by:
\begin{equation}
    \begin{split}
        C = & \sum_{\kvec, \mu} E_{\kvec, \mu} \frac{\partial n_{\mu}(\kvec)}{\partial T} 
        = \frac{\beta}{T} \sum_{\kvec, \mu} \frac{E_{\kvec, \mu}^2 e^{\beta E_{\kvec, \mu}}}{\left( e^{\beta E_{\kvec, \mu}} - 1 \right)^2}
        = k_{\rm B} \sum_{\kvec, \mu} \frac{(\beta E_{\kvec, \mu})^2 e^{\beta E_{\kvec, \mu}}}{\left( e^{\beta E_{\kvec, \mu}} - 1 \right)^2}
        =  k_{\rm B} \sum_{\kvec, \mu} \frac{ (\beta E_{\kvec, \mu})^2 }{\left( e^{\beta E_{\kvec, \mu}/2} - e^{-\beta E_{\kvec, \mu}/2} \right)^2} \\
        = &  k_{\rm B} \sum_{\kvec, \mu} \frac{ (\beta E_{\kvec, \mu})^2 }{4 \sinh^{2} (\beta E_{\kvec, \mu}/2 )}
    \end{split}
\end{equation}
This expression represents the specific heat contribution from magnon excitations in the system, which dominates the low-temperature thermal properties of magnetic insulators.



\subsection{Number and Spin moment from Correlation Matrix}\label{subsec: Number expectation value/Spin moment}

Finally, we discuss the number expectation and spin moment. 
Linear spin-wave theory gives the leading correction of the sublattice magnetization. 
The spin moments are 
\begin{equation}
    \langle \hat{\bf S}_{\mu} \rangle 
    \equiv \frac{1}{L} \sum_{i} \langle \hat{\bf S}_{i, \mu} \rangle
\end{equation}
where the spin moment at site $i=(i,\mu)$ can be obtained from $\langle \hat{\bf S}_{i, \mu} \rangle = \mathbf{R}_{i,\mu} \langle \widetilde{\bf S}_{i,\mu} \rangle$ and $\widetilde{\bf S}_{i,\mu} = (0, 0, \widetilde{S}^{0}_{i,\mu})$. 
This is the spin moment in the reference frame.
The average sublattice magnetization along the magnetization axis $\widetilde{S}^{0}_{i,\mu}$ is given
\begin{align*}
    \frac{1}{L} \sum_{i} \langle \widetilde{S}^{0}_{i, \mu} \rangle 
    = & \frac{1}{L} \sum_{i} \langle S_{\mu} -  \hat{n}_{i, \mu} \rangle 
    = S_{\mu} - \frac{1}{L} \sum_{\kvec \in \mathrm{FBZ}}  \langle \hat{n}_{\kvec, \mu}  \rangle
    = S_{\mu} - n_{\mu},
\end{align*}
where $i$ runs over the lattice, $\mu$ denotes the sublattice index, and $S_{\mu}$ is the spin moment of the spin. 
For example, for the spin 1/2 system, $S_\mu = 1/2$.
\begin{align}
    n_{\mu} \equiv & \frac{1}{L} \sum_{\kvec}  \langle \hat{n}_{\kvec, \mu} \rangle 
    = \frac{1}{L} \sum_{\kvec \in \mathrm{FBZ}} 
    = \langle b^{\dagger}_{\kvec, \mu} b_{\kvec, \mu}^{\phantom \dagger} \rangle  
\end{align}

To proceed, it is useful to introduce the correlation matrix.
The correlation matrix is defined as the expectation of outer product of the BdG vector $\Psi_{\kvec}$.
\begin{equation}
    \langle b^{\dagger}_{\kvec, \mu} b_{\kvec, \mu}^{\phantom \dagger} \rangle  
    = \Big[ \big\langle \Psi_{\kvec}^{\phantom \dagger} \Psi_{\kvec}^{\dagger} \big\rangle \Big]_{m_{s} + \mu, m_{s} + \mu} 
    = \left[
        \left\langle 
        \begin{pmatrix}
            {\bm b}_{\kvec} \\ {\bm b}^{\dagger}_{- \kvec}
        \end{pmatrix}
        \begin{pmatrix}
            {\bm b}_{\kvec}^{\dagger} & {\bm b}_{- \kvec}
        \end{pmatrix}
        \right\rangle
    \right]_{m_{s} + \mu, m_{s} + \mu} 
    = \left[
        \begin{pmatrix}
            \langle  {\bm b}_{\kvec}^{\phantom \dagger} {\bm b}_{\kvec}^{\dagger}  \rangle
            & \langle  {\bm b}_{\kvec}^{\phantom \dagger} {\bm b}_{- \kvec}^{\phantom \dagger}  \rangle
            \\ 
            \langle  {\bm b}^{\dagger}_{- \kvec} {\bm b}_{\kvec}^{\dagger}  \rangle
            & \langle  {\bm b}^{\dagger}_{- \kvec} {\bm b}_{- \kvec}^{\phantom \dagger}  \rangle
        \end{pmatrix}
    \right]_{m_{s} + \mu, m_{s} + \mu}
\end{equation}
where $\langle \hat{\bm b}_{\pm \kvec}^{(\dagger)} \hat{\bm b}_{\pm \kvec}^{(\dagger)} \rangle $ denotes the block matrix whose elements is the expectation for two bosonic operator ($\hat{\bm b}_{\pm \kvec}^{(\dagger)} \hat{\bm b}_{\pm \kvec}^{(\dagger)}$).
The expectation value is calculated for the state of the interest and our main interests are the ground states and the Gibbs states. 


The bosonic operator $(\hat{b}_{\kvec\sigma}^{\phantom \dagger}, \hat{b}_{\kvec\sigma}^{\dagger} )$ at $\kvec$ are related to the eigenmodes $(\hat{\beta}_{\kvec\sigma}^{\phantom \dagger}, \hat{\beta}_{\kvec\sigma}^{\dagger} )$ at momentum $\kvec$ through the para-unitary $\mathsf{T}_{\kvec}$. 
\begin{equation}
    \Psi_{\kvec} = \mathsf{T}_{\kvec} \widetilde{\Psi}_{\kvec}: \qquad
    \begin{pmatrix}
        \hat{\bm b}_{\kvec} \\ 
        \hat{\bm b}^{\dagger}_{-\kvec}
    \end{pmatrix}
    = 
    \begin{pmatrix}
        \mathsf{P}_{\kvec}           & \mathsf{Q}_{-\kvec}        \\
        \mathsf{Q}_{\kvec}^{*}       & \mathsf{P}_{-\kvec}^{*}
    \end{pmatrix}
    \begin{pmatrix}
        \hat{\bm\beta}_{\kvec} \\
        \hat{\bm\beta}_{-\kvec}^{\dagger}
    \end{pmatrix},
\end{equation}
where $\hat{\bm\beta}_{\kvec}^{T} = (\hat{\beta}_{\kvec,a}, \hat{\beta}_{\kvec,b}, \cdots)$.
Then, the correlation matrix can be obtained using unitary transformation.
\begin{equation}
    \big 
    \langle \Psi_{\kvec}^{\phantom \dagger} \Psi_{\kvec}^{\dagger} 
    \big \rangle 
    = 
    \big\langle 
    \mathsf{T}_{\kvec}^{\phantom\dagger} \widetilde{\Psi}_{\kvec}^{\phantom\dagger} \widetilde{\Psi}_{\kvec}^{\dagger} \mathsf{T}_{\kvec}^{\dagger} 
    \big\rangle
    = 
    \mathsf{T}_{\kvec} \big\langle \widetilde{\Psi}_{\kvec}^{\phantom\dagger} \widetilde{\Psi}_{\kvec}^{\dagger}  \big \rangle \mathsf{T}_{\kvec}^{\dagger}
    = 
    \mathsf{T}_{\kvec} 
    \begin{pmatrix}
            \langle \hat{\bm \beta}_{ \kvec}^{\phantom \dagger} \hat{\bm \beta}_{ \kvec}^{\dagger}          \rangle            
        &   \langle \hat{\bm \beta}_{ \kvec}^{\phantom \dagger} \hat{\bm \beta}_{-\kvec}^{\phantom \dagger} \rangle\\
            \langle \hat{\bm \beta}_{ \kvec}^{\dagger}          \hat{\bm \beta}_{-\kvec}^{\dagger}          \rangle  
        &   \langle \hat{\bm \beta}_{-\kvec}^{\dagger}          \hat{\bm \beta}_{-\kvec}^{\phantom \dagger} \rangle
    \end{pmatrix}
    \mathsf{T}_{\kvec}^{\dagger}
    =\mathsf{T}_{\kvec}^{\phantom \dagger} \mathsf{N}_{\kvec}^{\phantom \dagger} \mathsf{T}_{\kvec}^{\dagger},
\end{equation}
where $\mathsf{N}_{\kvec}$ is the $2m_{s} \times 2m_{s}$ Block matrix, 
For the Gibbs states, all elements for the correlation matrices can be obtained from the partition function. 
\begin{equation}
    \begin{array}{rlrl}
        \big[ \langle \hat{\bm \beta}_{ \kvec}^{\phantom \dagger} \hat{\bm \beta}_{ \kvec}^{\dagger} \rangle \big]_{\mu\nu} 
        & = \delta_{\mu\nu} \big( 1 + n_{\mu}(\kvec) \big) , 
        \qquad & \qquad 
        \big[ \langle \hat{\bm\beta}_{-\kvec}^{\dagger} \hat{\bm\beta}_{\kvec}^{\dagger} \rangle \big]_{\mu\nu} & = 0  ,
        \\[0.5em]
        \big[ \langle \hat{\bm\beta}_{-\kvec}^{\dagger} \hat{\bm\beta}_{- \kvec}^{\phantom \dagger} \rangle \big]_{\mu\nu}
        & = \delta_{\mu\nu} n_{\mu}(-\kvec) ,
        \qquad & \qquad 
        \big[ \langle \hat{\bm\beta}_{\kvec} \hat{\bm\beta}_{-\kvec} \rangle \big]_{\mu\nu} & = 0.
    \end{array}
\end{equation}
where $n_{\mu}(\kvec) = 1/(e^{\beta E_{\kvec}} - 1)$ is the Bose-Einstein distribution at energy $E_{\kvec}$. 








\newpage
\section{Correlations in Linear Spin Wave Theory}

In this section, we discuss the physical quantities related to spin-spin correlation functions. Specifically, from real-time (or dynamical) spin-spin correlation function to structure factor and spectral functions, these are quantities of fundamental interest in understanding magnetic excitations.

\begin{itemize}
    \item \textbf{Real-time correlation function:}
    The real-time spin-spin correlation function is a $3 \times 3$ matrix as a function of momentum and time:
    \begin{equation}
        \mathcal{C}^{\alpha\beta} (\mathbf{k}, t) = \frac{1}{L} \sum_{i, j} e^{- i \mathbf{k} \cdot (\mathbf{r}_{i} - \mathbf{r}_{j})} \langle \hat{S}_{i}^{\alpha}(t) \hat{S}_{j}^{\beta}(0) \rangle,
    \end{equation}
    where $\alpha, \beta = x, y, z$ or $\pm, 0$ and $L$ is the number of lattice sites for spins. 

    \begin{itemize}
        \item Note that the real-time spin-spin correlation function is equivalently defined as 
        \begin{equation}
            \mathcal{C}^{\alpha\beta} (\mathbf{k}, t) 
            = \big\langle \hat{S}_{+ \mathbf{k}}^{\alpha}(t) \hat{S}_{- \mathbf{k}}^{\beta}(0) \big\rangle,  
        \end{equation}
        by employing the spin Fourier transformation $S_{\mathbf{k}}^{\alpha} = \frac{1}{\sqrt{N}} \sum_{i} e^{-i \mathbf{k} \cdot \mathbf{r}_{i}} S_{i}^{\alpha}$.
        
        \item The spin-spin correlation function under complex conjugation is 
        \begin{equation}
            [\mathcal{C}^{\alpha\beta}(\mathbf{k}, t)]^{*}
            = \langle [S_{\mathbf{k}}^{\alpha} ( t ) S_{-\mathbf{k}}^{\beta} ( 0 )]^{\dagger} \rangle
            = \langle (S_{-\mathbf{k}}^{\beta} ( 0 ))^{\dagger} (S_{\mathbf{k}}^{\alpha} ( t ))^{\dagger} \rangle.
        \end{equation}
        At this point, the spin-spin correlation transforms differently depending on indices $\alpha, \beta$:
        
        \begin{equation}
            \langle (S_{-\mathbf{k}}^{\beta} ( 0 ))^{\dagger} (S_{\mathbf{k}}^{\alpha} ( t ))^{\dagger} \rangle
            = 
            \begin{cases}
                \langle S_{\mathbf{k}}^{\bar{\beta}} (0) S_{-\mathbf{k}}^{\bar{\alpha}} (t) \rangle
                = \langle S_{\mathbf{k}}^{\bar{\beta}} (-t) S_{-\mathbf{k}}^{\bar{\alpha}} (0) \rangle
                = \mathcal{C}^{\bar{\beta}\bar{\alpha}}(-\mathbf{k}, -t), 
                & \text{if } \alpha,\beta = \pm, 0, \\[1em]
                \langle S_{\mathbf{k}}^{\beta} (0) S_{-\mathbf{k}}^{\alpha} (t) \rangle
                = \langle S_{\mathbf{k}}^{\beta} (-t) S_{-\mathbf{k}}^{\alpha} (0) \rangle 
                = \mathcal{C}^{\beta\alpha}(-\mathbf{k}, -t), 
                & \text{if } \alpha, \beta = x, y, z.
            \end{cases}
        \end{equation}
        where in the second equality we use the identity $\langle \hat{A}(0) \hat{B}(t) \rangle = \langle \hat{A}(-t) \hat{B}(0) \rangle$ for thermal states $\rho = e^{-\beta H}/Z$:
        \begin{equation}
            \langle \hat{A}(0) \hat{B}(t) \rangle
            = \Tr[\rho \hat{A} \hat{U}^{\dagger}_{t} \hat{B} \hat{U}_{t}] 
            = \Tr[\rho \hat{U}_{t} \hat{A} \hat{U}^{\dagger}_{t} \hat{B}(0)]  
            = \langle \hat{A}(-t) \hat{B}(0) \rangle.
        \end{equation}
    \end{itemize}
    
    \item \textbf{Static structure factor:}
    The static structure factor is defined as
    \begin{equation}
        \mathcal{S}^{\alpha\beta}(\mathbf{k}) 
        \equiv
        \frac{1}{N} \sum_{i,j} e^{-i \mathbf{k} \cdot (\mathbf{r}_{i} - \mathbf{r}_{j})} 
        \big\langle \hat{S}^{\alpha}_{i} \hat{S}^{\beta}_{j} \big\rangle
        = \big\langle \hat{S}_{+ \mathbf{k}}^{\alpha} \hat{S}_{- \mathbf{k}}^{\beta} \big\rangle 
        = \mathcal{C}^{\alpha\beta}(\mathbf{k}, 0).
    \end{equation}
    This equals the equal-time spin-spin correlation function in momentum space, and can also be viewed as its time-average.

    
    \item \textbf{Dynamical structure factor:}
    The dynamic structure factor is defined as 
    \begin{equation}
        \mathcal{S}^{\alpha\beta}(\mathbf{k}, \omega) 
        \equiv \frac{1}{2\pi} \int_{-\infty}^{\infty} dt \, \mathcal{C}^{\alpha\beta}(\mathbf{k}, t) e^{i \omega t}.
    \end{equation}
    This quantity is directly related to the scattering cross-section in neutron scattering experiments.


    \item \textbf{Real space correlation function in the reference frame:}
    The spin-spin correlation in the local reference frame is:
    \begin{equation}
        \langle \widetilde{S}^{\alpha}_{i} \widetilde{S}_{j}^{\beta} \rangle,
    \end{equation}
    where $\widetilde{S}^{\alpha}_{i}$ represents the spin component in the local frame.
    
    \item \textbf{Spectral function:}
    The spectral function is defined as:
    \begin{equation}
        \mathcal{A}^{\alpha\beta}(\mathbf{k}, \omega) = -\frac{1}{\pi} \mathrm{Im}[G^{\alpha\beta}_{\rm R}(\mathbf{k}, \omega)],
    \end{equation}
    where $G^{\alpha\beta}_{\rm R}(\mathbf{k}, \omega)$ is the retarded Green's function:
    \begin{equation}
        G^{\alpha\beta}_{\rm R}(\mathbf{k}, \omega) 
        \equiv \frac{1}{i} \int_{-\infty}^{\infty} dt \, e^{i \omega t} \theta(t) \big\langle [S_{\mathbf{k}}^{\alpha}(t), S_{-\mathbf{k}}^{\beta}(0)] \big\rangle.
    \end{equation}
    The spectral function provides information about the density of states of the magnetic excitations.
\end{itemize}




\subsection{Real-Time (Dynamical) Spin-Spin Correlation Function}\label{subsec: Real-Time (Dynamical) Spin-Spin Correlation Function}

In this section, we derive the spin-spin correlation function in linear spin-wave theory (LSWT). The real-time spin-spin correlation function is a $3 \times 3$ matrix as a function of momentum and time:
\begin{equation}
    \mathcal{C}^{\alpha\beta}(\mathbf{k}, t) = \frac{1}{L} \sum_{i, j} e^{-i \mathbf{k} \cdot (\mathbf{r}_{i} - \mathbf{r}_{j})} \langle \hat{S}_{i}^{\alpha}(t) \hat{S}_{j}^{\beta}(0) \rangle,
\end{equation}
where $\alpha, \beta = x, y, z$ or $\pm, 0$ and $L$ is the number of lattice sites for spins. To proceed, we express the $\alpha$-component of the spin operator as 
\begin{equation}
    \hat{S}^{\alpha}_{i} 
    = \hat{\mathbf{e}}^{\alpha} \cdot \left( \frac{\widetilde{S}_{i}^{+} \hat{\mathbf{e}}^{-}_{i} + \widetilde{S}_{i}^{-} \hat{\mathbf{e}}^{+}_{i}}{\sqrt{2}} + \widetilde{S}^{0}_{i} \hat{\mathbf{e}}^{0}_{i} \right) 
    = \frac{1}{\sqrt{2}} \left( \bar{u}^{\alpha}_{i} \widetilde{S}_{i}^{+} + u^{\alpha}_{i} \widetilde{S}_{i}^{-} \right) + v^{\alpha}_{i} \widetilde{S}^{0}_{i},
\end{equation}
where $u^{\alpha}_{i} = \hat{\mathbf{e}}^{\alpha} \cdot \hat{\mathbf{e}}^{+}_{i}$, $\bar{u}^{\alpha}_{i} = (u^{\alpha}_{i})^*$, and $v^{\alpha}_{i} = \hat{\mathbf{e}}^{\alpha} \cdot \hat{\mathbf{e}}^{0}_{i}$ are components of the spin basis vectors. 
When the index $\alpha$ are $\pm,0$, then the $(u^{\alpha}_{i}, \bar{u}^{\alpha}_{i}, v^{\alpha}_{i})$ can be obtained from $\mathbf{Q}$:
\begin{equation}
    [\mathbf{Q}]_{\alpha\beta} = 
    \hat{\bf e}^{\alpha} \cdot \hat{\bf e}^{\beta}_{i} = 
    \hat{\bf e}^{\alpha} \cdot \mathbf{R}_{i} \hat{\bf e}^{\beta}: 
    \qquad
    \mathbf{Q}
    =
    \begin{pmatrix}
            u^{-}_{i} &   \bar{u}^{-}_{i}     &   v_{i}^{-}   \\
            u^{+}_{i} &   \bar{u}^{+}_{i}     &   v_{i}^{+}   \\
            u^{0}_{i} &   \bar{u}^{0}_{i}     &   v_{i}^{0}   \\
    \end{pmatrix}
    = 
    \begin{pmatrix}
        (\hat{\bf e}^{-})^{\rm T} \\  
        (\hat{\bf e}^{+})^{\rm T} \\  
        (\hat{\bf e}^{0})^{\rm T} 
    \end{pmatrix}
    \mathbf{R}_{i}
    \begin{pmatrix}
        \hat{\bf e}^{+} & \hat{\bf e}^{-} & \hat{\bf e}^{0}
    \end{pmatrix}
    = \mathbf{C}^{\dagger} \mathbf{R}_{i} \mathbf{C},
\end{equation}
If the indices $\alpha$ are $x,y,z$, $\mathbf{U}_{i} = \mathbf{R}_{i} \mathbf{C}$.
Substituting into the spin-spin correlation function, we obtain
\begin{align}
    \langle S_i^{\alpha}(t) S_j^{\beta}(0) \rangle
    &= \frac{1}{2} \Bigg\langle 
        \big( \bar{u}^{\alpha}_{i} \widetilde{S}_{i}^{+}(t) + u^{\alpha}_{i} \widetilde{S}_{i}^{-}(t) \big) 
        \big( \bar{u}^{\beta}_{j} \widetilde{S}_{j}^{+}(0) + u^{\beta}_{j} \widetilde{S}_{j}^{-}(0) \big) 
    \Bigg\rangle 
    + v^{\alpha}_{i} v^{\beta}_{j} \langle \widetilde{S}^{0}_{i}(t) \widetilde{S}^{0}_{j}(0) \rangle \nonumber \\
    &\quad + \frac{1}{\sqrt{2}} \Bigg[ 
        v^{\alpha}_{i} \Bigg\langle 
        \widetilde{S}^{0}_{i}(t) 
        \big( \bar{u}^{\beta}_{j} \widetilde{S}_{j}^{+}(0) + u^{\beta}_{j} \widetilde{S}_{j}^{-}(0) \big) 
        \Bigg\rangle 
        + v^{\beta}_{j} \Bigg\langle 
        \big( \bar{u}^{\alpha}_{i} \widetilde{S}_{i}^{+}(t) + u^{\alpha}_{i} \widetilde{S}_{i}^{-}(t) \big) 
        \widetilde{S}^{0}_{j}(0) 
        \Bigg\rangle 
    \Bigg].
\end{align}
In LSWT, which is a quadratic bosonic system, odd-order bosonic correlation functions vanish, so the terms in the second line are zero. Thus,
\begin{align}
    \langle S_i^{\alpha}(t) S_j^{\beta}(0) \rangle
    &= \frac{1}{2} \Bigg\langle 
        \big( \bar{u}^{\alpha}_{i} \widetilde{S}_{i}^{+}(t) + u^{\alpha}_{i} \widetilde{S}_{i}^{-}(t) \big) 
        \big( \bar{u}^{\beta}_{j} \widetilde{S}_{j}^{+}(0) + u^{\beta}_{j} \widetilde{S}_{j}^{-}(0) \big) 
    \Bigg\rangle 
    + v^{\alpha}_{i} v^{\beta}_{j} \langle \widetilde{S}^{0}_{i}(t) \widetilde{S}^{0}_{j}(0) \rangle \\
    &= \frac{1}{2} \sum_{m,n} \left[ 
    \begin{pmatrix}
        \bar{u}^{\alpha}_{i} & 0 \\ 0 & u^{\alpha}_{i}
    \end{pmatrix}
    \begin{pmatrix}
        \langle \widetilde{S}^{+}_{i}(t) \widetilde{S}^{-}_{j}(0) \rangle &
        \langle \widetilde{S}^{+}_{i}(t) \widetilde{S}^{+}_{j}(0) \rangle \\[0.25em]
        \langle \widetilde{S}^{-}_{i}(t) \widetilde{S}^{-}_{j}(0) \rangle &
        \langle \widetilde{S}^{-}_{i}(t) \widetilde{S}^{+}_{j}(0) \rangle
    \end{pmatrix}
    \begin{pmatrix}
        u^{\beta}_{j} & 0 \\ 0 & \bar{u}^{\beta}_{j}
    \end{pmatrix}
    \right]_{m,n}
    + v^{\alpha}_{i} v^{\beta}_{j} \langle \widetilde{S}^{0}_{i}(t) \widetilde{S}^{0}_{j}(0) \rangle.
\end{align}

Now, we define two matrix quantities: $\mathsf{U}^{\alpha}_{i}$, which maps the spin operators from the reference frame to the laboratory frame, and the correlation matrix between spin ladder operators (or transversal fluctuations).
\begin{equation}
    \mathsf{U}^{\alpha}_{i} 
    \equiv 
    \begin{pmatrix}
         \bar{u}^{\alpha}_{i} & 0 \\ 0 & u^{\alpha}_{i} 
    \end{pmatrix},
    \qquad \text{and} \qquad
    \big\langle \widetilde{\mathbf{S}}^{\pm}_{i}(t) \widetilde{\mathbf{S}}^{\mp}_{j}(0) \big\rangle
    = 
    \begin{pmatrix}
        \langle \widetilde{S}^{+}_{i}(t) \widetilde{S}^{-}_{j}(0) \rangle &
        \langle \widetilde{S}^{+}_{i}(t) \widetilde{S}^{+}_{j}(0) \rangle \\[0.25em]
        \langle \widetilde{S}^{-}_{i}(t) \widetilde{S}^{-}_{j}(0) \rangle &
        \langle \widetilde{S}^{-}_{i}(t) \widetilde{S}^{+}_{j}(0) \rangle
    \end{pmatrix}
\end{equation}
The Holstein-Primakoff (HP) operators are
\begin{equation}
    \widetilde{S}_{i}^{+} = \sqrt{2 S_i} \left( 1 - \frac{\hat{n}_i}{2 S_i} \right)^{1/2} \hat{b}_i, 
    \qquad \widetilde{S}_{i}^{-} = \sqrt{2 S_i} \hat{b}_i^{\dagger} \left( 1 - \frac{\hat{n}_i}{2 S_i} \right)^{1/2}, 
    \qquad \widetilde{S}_{i}^{0} = S_i - \hat{n}_i,
\end{equation}
with $\hat{n}_i = \hat{b}_i^{\dagger} \hat{b}_i$ being the magnon number operator.
We approximate $\widetilde{S}_{i}^{+} \approx \sqrt{2 S_i} \hat{b}_i$ and $\widetilde{S}_{i}^{-} \approx \sqrt{2 S_i} \hat{b}_i^{\dagger}$, yielding
\begin{equation}
    \langle S_i^{\alpha}(t) S_j^{\beta}(0) \rangle
    = \sqrt{S_i S_j} \sum_{m,n} \left[ 
    \mathsf{U}^{\alpha}_{i}
    \big\langle \Psi_{i}(t) \Psi_{j}^{\dagger}(0) \big\rangle
    \mathsf{U}^{\beta}_{j}
    \right]_{m,n}
    + v^{\alpha}_{i} v^{\beta}_{j} \big( S_i S_j -  S_i \langle \hat{n}_j(0) \rangle - S_j \langle \hat{n}_i(t) \rangle \big),
\end{equation}
where $\big\langle \Psi_{i}(t) \Psi_{j}^{\dagger}(0) \big\rangle$ is the matrix form of the two-point correlation function, defined as 
\begin{equation}
    \big\langle \Psi_{i}(t) \Psi_{j}^{\dagger}(0) \big\rangle
    \equiv
    \begin{pmatrix}
        \langle \hat{b}_i(t) \hat{b}_j^{\dagger}(0) \rangle &
        \langle \hat{b}_i(t) \hat{b}_j(0) \rangle \\
        \langle \hat{b}_i^{\dagger}(t) \hat{b}_j^{\dagger}(0) \rangle &
        \langle \hat{b}_i^{\dagger}(t) \hat{b}_j(0) \rangle
    \end{pmatrix}
    = \frac{1}{\sqrt{S_{i}S_{j}}}
    \big\langle \widetilde{\mathbf{S}}^{\pm}_{i}(t) \widetilde{\mathbf{S}}^{\mp}_{j}(0) \big\rangle
\end{equation}

The expectation value of the number operator satisfies $\langle \hat{n}_i(t) \rangle = \langle \hat{n}_i(0) \rangle$ because, for a time-independent Hamiltonian $\hat{H}$,
\begin{equation}
    \langle \hat{A}(t) \rangle = 
    \begin{cases}
        \langle \psi_n | e^{i \hat{H} t} \hat{A} e^{-i \hat{H} t} | \psi_n \rangle = \langle \hat{A} \rangle, & \text{if } \hat{H} | \psi_n \rangle = E_n | \psi_n \rangle, \\
        \mathrm{Tr} [ \rho \hat{U}^{\dagger}(t) \hat{A} \hat{U}(t) ] = \mathrm{Tr} [ \rho \hat{A} ], & \text{if } \rho = e^{-\beta \hat{H}} / Z,
    \end{cases}
\end{equation}
where $\hat{U}(t) = e^{-i \hat{H} t}$ is the time-evolution operator and $\rho = e^{-\beta \hat{H}} / Z$ is the thermal density matrix. In the Gibbs state, this follows from the commutation of $\rho$ with $\hat{U}(t)$ and the cyclic property of the trace.





\subsubsection{Sublattice Spin-Spin Correlation Function}

Given $L$ spins (or magnetic ions), there could be $m_{s}$ magnetic sublattices and $N$ unit cells with $L = m_{s}N$. 
In this case, we introduce a composite index notation: $I = (i, \mu)$ and $J = (j, \nu)$, where $\mu, \nu = a, b, \dots$ denote the magnetic sublattices and $i, j$ represent the unit cell indices. The position in the lattice is given by $\mathbf{R}_{I} = \mathbf{r}_{i} + \boldsymbol{\delta}_{\mu}$, where $\mathbf{r}_i$ is the position of the $i$-th unit cell and $\boldsymbol{\delta}_\mu$ is the position of the $\mu$-th sublattice within the unit cell.

Explicitly, for a spin at site $I = (i, \mu)$, its position vector is
\begin{equation}
    \mathbf{R}_{I} = \mathbf{R}_{i\mu} = \mathbf{r}_{i} + \boldsymbol{\delta}_{\mu}.
\end{equation}

To compute the real-time spin-spin correlation function in momentum space using Linear Spin-Wave Theory (LSWT), we express it in terms of magnetic sublattices:
\begin{align}
    \mathcal{C}^{\alpha\beta}(\mathbf{k}, t) 
    &= \frac{1}{Nm_{s}} \sum_{i\mu,j\nu} e^{-i\mathbf{k} \cdot [(\mathbf{r}_{i} + \boldsymbol{\delta}_{\mu}) - (\mathbf{r}_{j} + \boldsymbol{\delta}_{\nu})]} \langle \hat{S}_{i\mu}^{\alpha}(t) \hat{S}_{j\nu}^{\beta}(0) \rangle \\
    &= \frac{1}{m_{s}} \sum_{\mu,\nu} e^{-i\mathbf{k} \cdot (\boldsymbol{\delta}_{\mu} - \boldsymbol{\delta}_{\nu})} 
    \left[ 
    \frac{1}{N} \sum_{i,j} e^{-i\mathbf{k} \cdot (\mathbf{r}_{i} - \mathbf{r}_{j})} 
    \langle \hat{S}_{i\mu}^{\alpha}(t) \hat{S}_{j\nu}^{\beta}(0) \rangle \right] 
\end{align}

This allows us to express the total correlation function as an average over sublattice correlation functions:
\begin{align}
    \mathcal{C}^{\alpha\beta}(\mathbf{k}, t) 
    = \frac{1}{m_s} \sum_{\mu, \nu = a, b, \dots} 
    \mathcal{C}_{\mu\nu}^{\alpha\beta}(\mathbf{k}, t),
\end{align}
where $m_s$ is the number of magnetic sublattices in the unit cell, and $\mu, \nu$ denote sublattice indices.

The key insight of this section is to demonstrate that the sublattice real-time spin-spin correlation function $\mathcal{C}_{\mu\nu}^{\alpha\beta}(\mathbf{k}, t)$ can be obtained from a partial sum of the $2m_s \times 2m_s$ matrix $\mathsf{C}^{\alpha\beta}(\mathbf{k}, t)$:
\begin{equation}
    \begin{split}
        \mathcal{C}_{\mu\nu}^{\alpha\beta}(\mathbf{k}, t) 
        &= \sum_{(m,n)_{\mu,\nu}} 
        \big[ \mathsf{C}^{\alpha\beta}(\mathbf{k}, t) \big]_{m,n} \\
        &= \big[ \mathsf{C}^{\alpha\beta}(\mathbf{k}, t) \big]_{\mu, \nu} 
        + \big[ \mathsf{C}^{\alpha\beta}(\mathbf{k}, t) \big]_{m_s + \mu, \nu} 
        + \big[ \mathsf{C}^{\alpha\beta}(\mathbf{k}, t) \big]_{\mu, m_s + \nu} 
        + \big[ \mathsf{C}^{\alpha\beta}(\mathbf{k}, t) \big]_{m_s + \mu, m_s + \nu},
    \end{split}
\end{equation}
where the index set $(m,n)_{\mu,\nu} = \{ (\mu, \nu), (m_s + \mu, \nu), (\mu, m_s + \nu), (m_s + \mu, m_s + \nu) \}$ represents all possible combinations of indices in the extended Nambu space. The matrix $\mathsf{C}^{\alpha\beta}$ is defined as
\begin{equation}
    \boxed{
        \mathsf{C}^{\alpha\beta}(\mathbf{k}, t) 
        = 
        \mathsf{R}^{\alpha}_{\mathbf{k}} 
        \mathsf{T}_{\mathbf{k}} 
        \mathsf{N}_{\mathbf{k}}(t) 
        \mathsf{T}_{\mathbf{k}}^{\dagger}
        [\mathsf{R}^{\beta}_{\mathbf{k}}]^{\dagger},
    }
\end{equation}
where $\mathsf{T}_{\mathbf{k}}$ is a para-unitary matrix that diagonalizes the Hamiltonian matrix, satisfying $\mathsf{T}_{\mathbf{k}}^{\dagger} \mathsf{H}_{\mathbf{k}} \mathsf{T}_{\mathbf{k}} = \mathsf{E}_{\mathbf{k}}$, and $\mathsf{N}_{\mathbf{k}}(t)$ is the time-dependent correlation matrix in the magnon basis. 

The matrix $\mathsf{R}^{\alpha}_{\mathbf{k}}$ is a $2m_s \times 2m_s$ block-diagonal matrix that encodes both the spin components and sublattice structure. It is defined as $\mathsf{R}^{\alpha}_{\mathbf{k}} = \mathsf{U}^{\alpha} \mathsf{S}_{\mathbf{k}}$, which can be expressed explicitly as:
\begin{equation}
    \mathsf{R}^{\alpha}_{\mathbf{k}} 
    = \mathsf{U}^{\alpha} \mathsf{S}_{\mathbf{k}}
    =
    \begin{pmatrix}
        \bar{\mathsf{u}}^{\alpha} \mathsf{s}_{\mathbf{k}} & \mathbf{0} \\
        \mathbf{0} & \mathsf{u}^{\alpha} \mathsf{s}_{\mathbf{k}} 
    \end{pmatrix},
\end{equation}
where the diagonal matrices $\mathsf{s}$ and $\mathsf{u}^{\alpha}$ are defined as:
\begin{equation}
    \begin{cases}
        \mathsf{s} = \text{diag} \big[ e^{-i \mathbf{k} \cdot \boldsymbol{\delta}_a} \sqrt{S_a}, e^{-i \mathbf{k} \cdot \boldsymbol{\delta}_b} \sqrt{S_b}, \dots, e^{-i \mathbf{k} \cdot \boldsymbol{\delta}_\mu} \sqrt{S_\mu}, \dots \big], \\
        \mathsf{u}^{\alpha} = \text{diag} \big[ u_a^{\alpha}, u_b^{\alpha}, \dots, u_\mu^{\alpha}, \dots \big],
    \end{cases}
\end{equation}
where $\mathsf{u}^{\alpha}$ contains the $\alpha$-component of the spin basis vector $\hat{\mathbf{e}}^-$ for each sublattice, i.e., $\mathsf{u}^{\alpha}_\mu = [\hat{\mathbf{e}}^+_\mu]^{\alpha}$, and $\bar{\mathsf{u}}^{\alpha} = [\mathsf{u}^{\alpha}]^* = [\mathsf{u}^{\alpha}]^{\dagger}$. The matrix $\mathsf{s}$ encodes both the phase factors due to sublattice positions and the square root of spin magnitudes $S_\mu$.

The matrix $\mathsf{N}_{\mathbf{k}}(t)$ encodes the time evolution of magnon correlations and is given by:
\begin{equation}
    \mathsf{N}_{\mathbf{k}}(t) 
    = 
    \begin{pmatrix}
        \langle \hat{\boldsymbol{\beta}}_{\mathbf{k}}(t) \hat{\boldsymbol{\beta}}_{\mathbf{k}}^{\dagger}(0) \rangle & \langle \hat{\boldsymbol{\beta}}_{\mathbf{k}}(t) \hat{\boldsymbol{\beta}}_{-\mathbf{k}}(0) \rangle \\
        \langle \hat{\boldsymbol{\beta}}_{-\mathbf{k}}^{\dagger}(t) \hat{\boldsymbol{\beta}}_{\mathbf{k}}^{\dagger}(0) \rangle & \langle \hat{\boldsymbol{\beta}}_{-\mathbf{k}}^{\dagger}(t) \hat{\boldsymbol{\beta}}_{-\mathbf{k}}(0) \rangle
    \end{pmatrix}
    = 
    \begin{pmatrix}
        (\mathsf{I} + \mathsf{n}_{\mathbf{k}})e^{-it\mathsf{E}_{\mathbf{k}}} & \mathbf{0} \\
        \mathbf{0} & \mathsf{n}_{-\mathbf{k}}e^{it\mathsf{E}_{-\mathbf{k}}}
    \end{pmatrix},
\end{equation}
where $\mathsf{I}$ is the $m_s \times m_s$ identity matrix, $\mathsf{n}_{\mathbf{k}}$ and $\mathsf{E}_{\mathbf{k}}$ are diagonal matrices containing the Bose-Einstein distribution $[\mathsf{n}_{\mathbf{k}}]_{\mu\nu} = \delta_{\mu\nu}/[e^{\beta E_\mu(\mathbf{k})} - 1]$ and the magnon energies $[\mathsf{E}_{\mathbf{k}}]_{\mu\nu} = \delta_{\mu\nu}E_{\mu}(\mathbf{k})$, respectively. The time dependence arises from the energy eigenstates: $e^{i\hat{H}t}\hat{\beta}_{\mathbf{k},\mu}e^{-i\hat{H}t} = e^{-iE_{\mathbf{k},\mu}t}\hat{\beta}_{\mathbf{k},\mu}$.



To derive this, we assume a spin system with magnetic sublattice structures within the unit cell. The sublattice correlation function is expressed as
\begin{align}
    \mathcal{C}^{\alpha\beta}_{\mu\nu} (\kvec, t) 
    = & \frac{1}{N} \sum_{i,j} e^{ - i \kvec \cdot [(\mathbf{r}_{i} + \mathbf{\delta}_{\mu}) - (\mathbf{r}_{j} + \mathbf{\delta}_{\nu})]} \langle S_{i\mu}^{\alpha}(t) S_{j\nu}^{\beta}(0) \rangle \\
    = & e^{ - i \kvec \cdot (\mathbf{\delta}_{\mu} - \mathbf{\delta}_{\nu})}  
    \frac{1}{N} \sum_{i,j} e^{ - i \kvec \cdot (\mathbf{r}_{i} - \mathbf{r}_{j} )} 
    \langle S_{i\mu}^{\alpha}(t) S_{j\nu}^{\beta}(0) \rangle \\
    = & e^{ - i \kvec \cdot (\mathbf{\delta}_{\mu} - \mathbf{\delta}_{\nu})}  
    \frac{1}{N}  \\
    & \times \sum_{i,j} e^{ - i \kvec \cdot (\mathbf{r}_{i} - \mathbf{r}_{j} )} 
    \bigg[
    \sqrt{S_{\mu}S_{\nu}} \sum_{m,n}
    \Big[ \mathsf{U}^{\alpha}_{\mu}
    \big\langle \Psi_{i\mu}^{\phantom \dagger}(t) \Psi_{j\nu}^{\dagger}(0) \big\rangle
    \mathsf{U}^{\beta}_{\nu} \Big]_{m,n} 
    + v_{\mu}^{\alpha} v_{\nu}^{\beta}
    \Big( S_{\nu} S_{\mu} - S_{\nu}\langle \hat{n}_{i\mu} \rangle - S_{\mu}\langle \hat{n}_{j\nu} \rangle \Big)
    \bigg], 
    \nonumber
\end{align}
where the spin operators are transformed into bosonic operators via the Holstein-Primakoff transformation, and the correlation is computed using Fourier transforms. The first term, involving the bosonic correlation matrix, yields
\begin{subequations}
    \begin{align}
        \frac{1}{N} \sum_{i,j}   
        e^{ - i \kvec \cdot (\mathbf{r}_{i} - \mathbf{r}_{j} )} 
        \big\langle \Psi_{i}^{\phantom \dagger}(t) \Psi_{j}^{\dagger}(0) \big\rangle 
        = &
        \big\langle \Psi_{\kvec}^{\phantom \dagger}(t) \Psi_{\kvec}^{\dagger}(0) \big\rangle,
        \\
        \frac{1}{N} \sum_{\kvec} 
        e^{ + i \kvec \cdot (\mathbf{r}_{i} - \mathbf{r}_{j} )} 
        \big\langle \Psi_{\kvec}^{\phantom \dagger}(t) \Psi_{\kvec}^{\dagger}(0) \big\rangle 
        = &
        \big\langle \Psi_{i}^{\phantom \dagger}(t) \Psi_{j}^{\dagger}(0) \big\rangle.
    \end{align}
\end{subequations}
Similarly, the spin-spin correlation function in real space is obtained as
\begin{equation}
    \big\langle \widetilde{\mathbf{S}}_{i}^{\pm}(t) \widetilde{\mathbf{S}}_{j}^{\mp}(0) \big\rangle
    = \frac{1}{N} \sum_{\kvec, \mu, \nu } 
    e^{ + i \kvec \cdot (\mathbf{r}_{i} - \mathbf{r}_{j} )} 
    \sqrt{S_{\mu} S_{\nu}} 
    e^{ i \kvec \cdot (\mathbf{\delta}_{\mu} - \mathbf{\delta}_{\nu}) }
    \big\langle \Psi_{\kvec, \mu}^{\phantom \dagger}(t) \Psi_{\kvec, \nu}^{\dagger}(0) \big\rangle
    = \frac{1}{N} \sum_{\kvec} 
    e^{ + i \kvec \cdot (\mathbf{r}_{i} - \mathbf{r}_{j} )} 
    \bar{\mathsf{S}}_{\kvec}^{\phantom \dagger} 
    \mathsf{T}_{\kvec}^{\phantom \dagger}
    \mathsf{N}_{\kvec}^{\phantom \dagger} (t)
    \mathsf{T}_{\kvec}^{\dagger} 
    \mathsf{S}_{\kvec}^{\phantom \dagger},
\end{equation}
derived by applying the Fourier transform to the bosonic operators and computing the correlation matrix elements:
\begin{align}
    \frac{1}{N} \sum_{i,j} e^{ - i \kvec \cdot (\mathbf{r}_{i} - \mathbf{r}_{j} )} 
    \big\langle \Psi_{i\mu}^{\phantom \dagger}(t) \Psi_{j\nu}^{\dagger}(0) \big\rangle 
    = & 
    \begin{pmatrix}
        \big\langle \hat{b}_{+\kvec, \mu}^{\phantom \dagger} (t) \hat{b}_{+\kvec, \nu}^{\dagger}          (0) \big\rangle &
        \big\langle \hat{b}_{+\kvec, \mu}^{\phantom \dagger} (t) \hat{b}_{-\kvec, \nu}^{\phantom \dagger} (0) \big\rangle \\
        \big\langle \hat{b}_{-\kvec, \mu}^{\dagger}          (t) \hat{b}_{+\kvec, \nu}^{\dagger}          (0) \big\rangle &
        \big\langle \hat{b}_{-\kvec, \mu}^{\phantom \dagger} (t) \hat{b}_{-\kvec, \nu}^{\dagger}          (0) \big\rangle 
    \end{pmatrix}
    = \mathsf{T}_{\kvec}^{\phantom\dagger} \mathsf{N}_{\kvec}^{\phantom\dagger} (t) \mathsf{T}_{\kvec}^{\dagger},
\end{align}
where $\mathsf{T}_{\kvec}$ and $\mathsf{N}_{\kvec}(t)$ are defined above. The second term, representing the density operator in momentum space, is
\begin{equation}
    \frac{1}{N} \sum_{i,j} e^{- i \kvec \cdot (\mathbf{r}_{i} - \mathbf{r}_{j})} \langle \hat{n}_{j\nu} \rangle
    = \Big( \sum_{\kappa \in \mathbf{G}} \delta(\kvec - \kappa) \Big) \langle \hat{n}_{\kvec, \mu} \rangle,
\end{equation}
where $\hat{n}_{\kvec, \mu}$, the density fluctuation operator at momentum $\kvec$, is
\begin{align*}
    \hat{n}_{\kvec} 
    = & \frac{1}{N} 
    \sum_{i} e^{i \kvec \cdot \mathbf{r}_{i}} \hat{a}^{\dagger}_{i} \hat{a}_{i}^{\phantom \dagger} 
    = \frac{1}{N}  \sum_{\mathbf{q}} \hat{a}_{\mathbf{q} + \kvec}^{\dagger} \hat{a}_{\mathbf{q}}^{\phantom \dagger},
\end{align*}
which vanishes for a free system unless $\kvec = \kappa \in \mathbf{G}$ due to momentum conservation.

The correlation function is expressed in matrix form as
\begin{equation}
    \mathcal{C}^{\alpha\beta}_{\mu\nu} (\kvec, t)
    = e^{ - i \kvec \cdot (\mathbf{\delta}_{\mu} - \mathbf{\delta}_{\nu})}  
    \left[
    \sqrt{S_{\mu}S_{\nu}} \sum_{m,n}
    \Big[ \mathsf{U}^{\alpha}_{\mu}
    \mathsf{T}_{\kvec}^{\phantom\dagger} \mathsf{N}_{\kvec}^{\phantom\dagger} (t) \mathsf{T}_{\kvec}^{\dagger}
    \mathsf{U}^{\beta}_{\nu} \Big]_{m,n}
    + v_{\mu}^{\alpha} v_{\nu}^{\beta}
    \Big( \sum_{\kappa \in \mathbf{G}} \delta(\kvec - \kappa) \Big)
    \Big( S_{\nu}\langle \hat{n}_{\mathbf{0}, \mu} \rangle + S_{\mu}\langle \hat{n}_{\mathbf{0}, \nu} \rangle \Big)
    \right],
\end{equation}
where the first term describes time-dependent magnon scattering, and the second term accounts for the static ordered moment reduction due to magnon population. The matrix form is
\begin{equation}\label{eq: matrix form of spin-spin correlation matrix}
    \mathsf{C}^{\alpha\beta} (\kvec, t)
    = 
    \mathsf{U}^{\alpha} 
    \mathsf{S}_{\kvec}^{\phantom \dagger} 
    \mathsf{T}_{\kvec}^{\phantom \dagger} 
    \mathsf{N}_{\kvec}^{\phantom \dagger}  (t)
    \mathsf{T}_{\kvec}^{         \dagger}
    \mathsf{S}_{\kvec}^{         \dagger}
    \mathsf{U}^{\beta},
\end{equation}
with elements computed as
\begin{align*}
    \mathsf{C}^{\alpha\beta} (\kvec, t) = & 
    \begin{pmatrix}
        \bar{\mathsf{u}}^{\alpha} \mathsf{s}_{\kvec}  & \mathbf{0} \\
        \mathbf{0} & \mathsf{u}^{\alpha} \mathsf{s}_{\kvec} \\
    \end{pmatrix}
    \begin{pmatrix}     
        \langle \mathbf{b}_{ \mathbf{k}\mu}           (t)    \mathbf{b}_{ \mathbf{k}\nu}^{\dagger} (0)    \rangle & 
        \langle \mathbf{b}_{ \mathbf{k}\mu}           (t)    \mathbf{b}_{-\mathbf{k}\nu}           (0)    \rangle \\
        \langle \mathbf{b}_{-\mathbf{k}\mu}^{\dagger} (t)    \mathbf{b}_{ \mathbf{k}\nu}^{\dagger} (0)    \rangle & 
        \langle \mathbf{b}_{-\mathbf{k}\mu}^{\dagger} (t)    \mathbf{b}_{-\mathbf{k}\nu}           (0)    \rangle
    \end{pmatrix}
    \begin{pmatrix}
        \bar{\mathsf{s}}_{\kvec} \mathsf{u}^{\beta} & \mathbf{0} \\
        \mathbf{0} & \bar{\mathsf{s}}_{\kvec} \bar{\mathsf{u}}^{\beta} \\
    \end{pmatrix}.
\end{align*}
The elements of $\mathsf{C}^{\alpha\beta}(\kvec, t)$ for indices $(m,n)_{\mu,\nu}$ are
\begin{align*}
    \begin{array}{crl}
        {[\mathsf{C}^{\alpha\beta}(\kvec, t)]_{(1,1)_{\mu\nu}}} : \qquad\qquad
        &
        \big[
            \mathsf{s}_{\kvec} \bar{\mathsf{u}}^{\alpha} 
            \langle \mathbf{b}_{+\mathbf{k}\mu}           (t)    \mathbf{b}_{+\mathbf{k}\nu}^{\dagger} (0)    \rangle 
            \mathsf{u}^{\beta} \bar{\mathsf{s}}_{\kvec} 
        \big]_{\mu\nu}
        = &
        \sqrt{S_{\mu}S_{\nu}}e^{-i \kvec \cdot (\bm{\delta}_{\mu} - \bm{\delta}_{\nu})} 
        \big\langle b_{+\mathbf{k}\mu} (t) b_{+\mathbf{k}\nu}^{\dagger} (0) \big\rangle,
        \\
        {[\mathsf{C}^{\alpha\beta}(\kvec, t)]_{(1,2)_{\mu\nu}}} : \qquad\qquad
        &
        \big[
            \mathsf{s}_{\kvec} \bar{\mathsf{u}}^{\alpha} 
            \langle \mathbf{b}_{+\mathbf{k}\mu}           (t)    \mathbf{b}_{-\mathbf{k}\nu}           (0)    \rangle 
            \bar{\mathsf{u}}^{\beta} \bar{\mathsf{s}}_{\kvec}
        \big]_{\mu\nu}
        = & \sqrt{S_{\mu}S_{\nu}}e^{-i \kvec \cdot (\bm{\delta}_{\mu} - \bm{\delta}_{\nu})} 
        \big\langle b_{+\mathbf{k}\mu} (t) b_{-\mathbf{k}\nu} (0) \big\rangle,
        \\
        {[\mathsf{C}^{\alpha\beta}(\kvec, t)]_{(2,1)_{\mu\nu}}} : \qquad\qquad
        & 
        \big[
            \mathsf{s}_{\kvec} \mathsf{u}^{\alpha}
            \langle \mathbf{b}_{-\mathbf{k}\mu}^{\dagger} (t)    \mathbf{b}_{+\mathbf{k}\nu}^{\dagger} (0)    \rangle  
            \mathsf{u}^{\beta} \bar{\mathsf{s}}_{\kvec}  
        \big]_{\mu\nu}
        = & \sqrt{S_{\mu}S_{\nu}}e^{-i \kvec \cdot (\bm{\delta}_{\mu} - \bm{\delta}_{\nu})} 
        \big\langle b_{-\mathbf{k}\mu}^{\dagger} (t) b_{+\mathbf{k}\nu}^{\dagger} (0) \big\rangle,
        \\
        {[\mathsf{C}^{\alpha\beta}(\kvec, t)]_{(2,2)_{\mu\nu}}} : \qquad\qquad
        &
        \big[
            \mathsf{s}_{\kvec} \mathsf{u}^{\alpha}
            \langle \mathbf{b}_{-\mathbf{k}\mu}^{\dagger} (t)    \mathbf{b}_{-\mathbf{k}\nu}           (0)    \rangle
            \bar{\mathsf{u}}^{\beta} \bar{\mathsf{s}}_{\kvec}
        \big]_{\mu\nu}
        = & \sqrt{S_{\mu}S_{\nu}}e^{-i \kvec \cdot (\bm{\delta}_{\mu} - \bm{\delta}_{\nu})} 
        \big\langle b_{-\mathbf{k}\mu}^{\dagger} (t) b_{-\mathbf{k}\nu} (0) \big\rangle.
    \end{array}
\end{align*}
Summing $\mathsf{C}^{\alpha\beta}(\kvec, t)$ over $(m,n)_{\mu,\nu}$ yields the sublattice spin-spin correlation function
\begin{equation}
    \mathcal{C}_{\mu\nu}^{\alpha\beta} (\kvec, t) = \sum_{(m,n)_{\mu,\nu}} [\mathsf{C}^{\alpha\beta}(\kvec, t)]_{m,n}.
\end{equation}
The spin structure factor is obtained by summing over all sublattices: $\mathcal{S}^{\alpha\beta} (\mathbf{k}, t) = \mathcal{C}^{\alpha\beta} (\kvec, t) = \frac{1}{m_s} \sum_{\mu,\nu} \mathcal{C}_{\mu\nu}^{\alpha\beta} (\kvec, t)$.




\subsection{Structure Factor}\label{subsec: Structure Factor}

In this section, we derive the formulas for the static and dynamic structure factors in Linear Spin-Wave Theory.

\subsubsection{Static Structure Factor}

The static structure factor in momentum space is given by
\begin{equation}
    \mathcal{S}^{\alpha\beta} (\mathbf{k}) 
    \equiv \frac{1}{L} 
    \sum_{I, J}
    e^{- i \mathbf{k} \cdot ( \mathbf{r}_{I} -  \mathbf{r}_{J} ) }
    \langle 
    S_{I}^{\alpha}
    S_{J}^{\beta} 
    \rangle 
    = \big\langle S^{\alpha}_{\kvec} S^{\beta}_{-\kvec} \big\rangle.
\end{equation}
This quantity $\mathcal{S}^{\alpha\beta} (\mathbf{k})$ is equal to $\mathcal{C}^{\alpha\beta}(\kvec, t)$, which we've found the matrix form in a previous equation:
\begin{equation}
    \mathcal{S}^{\alpha\beta} (\mathbf{k}) = 
    \sum_{m,n} 
    \big[ \mathsf{U}^{\alpha} 
    \mathsf{S}_{\kvec}^{\phantom \dagger} 
    \mathsf{T}_{\kvec}^{\phantom \dagger} 
    \mathsf{N}_{\kvec}^{\phantom \dagger}
    \mathsf{T}_{\kvec}^{         \dagger}
    \mathsf{S}_{\kvec}^{         \dagger}
    \mathsf{U}^{\beta} \big]_{mn},
\end{equation}
where $\mathsf{U}^{\alpha}$ and $\mathsf{U}^{\beta}$ are spin projection operators, $\mathsf{S}_{\kvec}$ and $\mathsf{T}_{\kvec}$ are transformation matrices, and $\mathsf{N}_{\kvec}$ accounts for the diagonalized magnon correlation functions.
In the reference frame, the structure factor is expressed as
\begin{align*}
    \widetilde{\mathcal{S}}^{\alpha\beta} (\mathbf{k}) 
    \equiv &
    \frac{1}{L} 
    \sum_{I, J}
    e^{- i \mathbf{k} \cdot ( \mathbf{r}_{I} -  \mathbf{r}_{J} ) }
    \langle \widetilde{S}_{I}^{\alpha} \widetilde{S}_{J}^{\beta} \rangle 
    = 
    \sum_{\mu, \nu} 
    \left[
    \begin{pmatrix}
        \langle \widetilde{S}^{+}_{\kvec} \widetilde{S}^{-}_{-\kvec} \rangle &
        \langle \widetilde{S}^{+}_{\kvec} \widetilde{S}^{+}_{-\kvec} \rangle \\[0.25em]
        \langle \widetilde{S}^{-}_{\kvec} \widetilde{S}^{-}_{-\kvec} \rangle &
        \langle \widetilde{S}^{-}_{\kvec} \widetilde{S}^{+}_{-\kvec} \rangle
    \end{pmatrix} 
    \right]_{(\alpha, \beta)_{\mu\nu}}
    \\ = &
    \sum_{\mu, \nu} 
    \big[
    \mathsf{S}_{\kvec}^{\phantom \dagger} 
    \mathsf{T}_{\kvec}^{\phantom \dagger} 
    \mathsf{N}_{\kvec}^{\phantom \dagger}  
    \mathsf{T}_{\kvec}^{         \dagger}
    \mathsf{S}_{\kvec}^{         \dagger}
    \big]_{(\alpha, \beta)_{\mu\nu}},
\end{align*}
where the matrix elements are evaluated in the rotated spin basis.

The quantity relevant to neutron scattering experiments is defined as:
\begin{equation}
    \mathcal{S} (\mathbf{q}) 
    \equiv 
    \frac{1}{N} 
    \sum_{\substack{\alpha, \beta \\ I, J}}
    \Big(
    \delta_{\alpha\beta} - \frac{q_{\alpha} q_{\beta}}{q^{2}} 
    \Big)
    e^{i \mathbf{q} \cdot ( \mathbf{r}_{I} - \mathbf{r}_{J} ) }
    \langle 
    S_{i}^{\alpha}
    S_{j}^{\beta} 
    \rangle
    = 
    \sum_{\alpha, \beta} 
    \Big( \delta_{\alpha\beta} - \frac{q_{\alpha} q_{\beta}}{q^{2}} \Big)
    \mathcal{S}^{\alpha\beta} (\mathbf{q}) ,
\end{equation}
This formulation respects the directional constraints intrinsic to neutron scattering physics.



\subsubsection{Dynamic structure factor} 

The dynamic structure factor is defined as 
\begin{equation}
    \mathcal{S}^{\alpha\beta}(\kvec, \omega) 
    = \frac{1}{2\pi} \int_{-\infty}^{\infty} d t \, e^{i \omega t} \mathcal{C}^{\alpha\beta} (\kvec, t).
\end{equation}
From the LSWT, we consider 
\begin{align*}
    \mathcal{S}^{\alpha\beta}_{\mu\nu} (\mathbf{k}, \omega) 
    = & \sum_{(m,n)_{\mu,\nu}} \int_{-\infty}^{\infty} dt \, [\mathsf{C}^{\alpha\beta} (\kvec, t)]_{m,n} e^{i \omega t}
    = \sum_{(m,n)_{\mu,\nu}} 
    \left[ 
    \int_{-\infty}^{\infty} dt \, e^{i \omega t} 
    \mathsf{U}^{\alpha} 
    \mathsf{S}_{\kvec}^{\phantom \dagger} 
    \mathsf{T}_{\kvec}^{\phantom \dagger} 
    \mathsf{N}_{\kvec}^{\phantom \dagger}  (t)
    \mathsf{T}_{\kvec}^{         \dagger}
    \mathsf{S}_{\kvec}^{         \dagger}
    \mathsf{U}^{\beta}
    \right]_{m,n}  \\
    = & \sum_{(m,n)_{\mu,\nu}} 
    \left[  \mathsf{R}^{\alpha}_{\kvec} \mathsf{T}_{\kvec}
    \left( \int_{-\infty}^{\infty} dt \, e^{i \omega t} \mathsf{N}_{\kvec} (t) \right)
    \mathsf{T}_{\kvec}^{\dagger} [\mathsf{R}^{\beta}_{\kvec}]^{\dagger} \right]_{m,n}.
\end{align*}
However, the integral yields $\delta$-function. So we introduce the decay factor $\eta > 0$ so that $\omega \rightarrow \omega + i \eta$:
\begin{equation}
    g_{\kvec} ( \omega ; \eta ) 
    \equiv \int_{0}^{\infty} dt \, e^{ - [ \eta  + i( E_{\kvec} - \omega ) ] t} 
    = \frac{1}{\eta + i ( E_{\kvec} - \omega )}
\end{equation}
Using this notation, we have
\begin{align}
    \mathcal{F}_{\kvec} (\omega; \eta) = & \int_{-\infty}^{\infty} dt\, e^{-i E_{\kvec} t} e^{ i\omega t - \eta |t|} 
    =  \int_{0}^{\infty} dt \, 
    \big( e^{ - [ \eta  + i( E_{\kvec} - \omega ) ] t} + e^{ - [ \eta  - i( E_{\kvec} - \omega ) ] t} \big) 
    =  g_{\kvec} ( \omega ; \eta ) + g_{\kvec}^{*} ( \omega ; \eta ) \\
    = & 2 \mathrm{Re} [ g_{\kvec} ( \omega ; \eta ) ]
    = \frac{2\eta}{\eta^{2} + ( E_{\kvec} - \omega )^{2}}
\end{align}
where $\mathcal{F}_{\kvec}(\omega; \eta)$ is Lorentzian and its FWHM is $2\eta$.
\begin{equation}
    \begin{split}
        \int_{-\infty}^{\infty} dt \, e^{i \omega t - \eta |t|} \mathsf{N}_{\kvec} (t) 
        = &
        \int_{-\infty}^{\infty} dt \, e^{i \omega t  - \eta |t|}
        \begin{pmatrix}
            ( \mathsf{I} + \mathsf{n}_{\kvec} ) e^{-{ i E}_{\kvec,\mu}t}  & 0 \\
            0                         & \mathsf{n}_{-\kvec}  e^{+{  i  E}_{-\kvec,\mu}t}
        \end{pmatrix} \\
        = &
        \begin{pmatrix}
            ( \mathsf{I} + \mathsf{n}_{\kvec} ) 
            \int_{-\infty}^{\infty} dt \, e^{i \omega t  - \eta |t|} 
            e^{-{ i E}_{\kvec,\mu}t}  & 0 \\
            0                         & \mathsf{n}_{-\kvec}  
            \int_{-\infty}^{\infty} dt \, e^{i \omega t  - \eta |t|}  e^{+{  i  E}_{-\kvec,\mu}t}
        \end{pmatrix} \\
        = & 
        \begin{pmatrix}
            ( \mathsf{I} + \mathsf{n}_{\kvec} ) \mathcal{F}_{\kvec} (\omega; \eta)  & 0 \\
            0   &   \mathsf{n}_{-\kvec} \mathcal{F}_{- \kvec} (-\omega; \eta)
        \end{pmatrix} 
    \end{split}
\end{equation}
Thus, in LSWT, the dynamical spin structure factor can be computed from the matrix multiplication:
\begin{equation}
    \mathcal{S}^{\alpha\beta}_{\mu\nu} (\kvec, \omega)
    =
    \sum_{(m,n)_{\mu,\nu}} 
    \left[  \mathsf{R}^{\alpha}_{\kvec} \mathsf{T}_{\kvec}
    \mathcal{F} [\mathsf{N}_{\kvec}](\omega; \eta)
    \mathsf{T}_{\kvec}^{\dagger} [\mathsf{R}^{\beta}_{\kvec}]^{\dagger} \right]_{m,n}.
\end{equation}
Here, the fourier transformation of diagonal matrix $\mathsf{N}_{\kvec} (t)$ is given by
\begin{equation}
    \mathcal{F} [\mathsf{N}_{\kvec}](\omega; \eta)
    \equiv 
    \begin{pmatrix}
        ( \mathsf{I} + \mathsf{n}_{\kvec} )    \mathcal{F}_{ \kvec} (  \omega; \eta  )  & 0 \\
        0   &   \mathsf{n}_{-\kvec} \mathcal{F}_{-\kvec} ( -\omega; \eta  )
    \end{pmatrix} 
\end{equation}


\subsection{Spectral Function}\label{subsec: Spectral Function}

In this section, we compute the spectral function within Linear Spin-Wave Theory (LSWT):
\begin{equation}
    \mathcal{A}^{\alpha\beta}(\kvec, \omega) = - \frac{1}{\pi} \mathrm{Im} [G^{\alpha\beta}_{\rm R} (\kvec, \omega)],
\end{equation}
where $G^{\alpha\beta}_{\rm R} (\kvec, \omega)$ is the retarded Green's function, defined as
\begin{equation}
    G^{\alpha\beta}_{\rm R} (\kvec, \omega) 
    \equiv  \frac{1}{i} \int_{-\infty}^{\infty} dt \, e^{i\omega t} \theta(t) \big\langle [S_{\kvec}^{\alpha} (t), S_{-\kvec}^{\beta} (0)] \big\rangle . 
\end{equation}
The ratared Green's function can be rewritten in terms of the correlation function
\begin{equation}
    G^{\alpha\beta}_{\rm R} (\kvec, \omega) 
    = \frac{1}{i} \int_{0}^{\infty} dt \, \Big[ \big\langle S_{\kvec}^{\alpha} (t) S_{-\kvec}^{\beta} (0) \big\rangle - \big\langle S_{-\kvec}^{\beta} (0) S_{\kvec}^{\alpha} (t) \big\rangle \Big] e^{i \omega t} 
    = \frac{1}{i} \int_{0}^{\infty} dt \, \big[ \mathcal{C}^{\alpha\beta}(\kvec, t) - \mathcal{C}^{\beta\alpha}(-\kvec, -t) \big] e^{i \omega t}.
\end{equation}
For spin components $\alpha, \beta = x, y, z$, the retarded Green's function is expressed as the Fourier transform of the imaginary part of the correlation function:
\begin{equation}
    G^{\alpha\beta}_{\rm R} (\kvec, \omega) 
    = 2 \int_{0}^{\infty} dt \, \mathrm{Im} [\mathcal{C}^{\alpha\beta}(\kvec, t) ] e^{\ri \omega t}.
\end{equation}
The spectral function is then obtained as
\begin{equation}
    \begin{split}
        \mathcal{A}^{\alpha\beta}(\kvec, \omega) 
        % = & - \frac{1}{\pi} \mathrm{Im} [G^{\alpha\beta}_{\rm R} (\kvec, \omega)] \\
        = & - \frac{2}{\pi} \int_{0}^{\infty} dt \, \mathrm{Im} \big[ \mathrm{Im} \big[ \mathcal{C}^{\alpha\beta}(\kvec, t) \big] e^{i \omega t} \big] \\
        = & - \frac{2}{\pi} \int_{0}^{\infty} dt \, \mathrm{Im} \big[ \mathcal{C}^{\alpha\beta}(\kvec, t) \sin (\omega t) \big] \\
        = & - \frac{2}{\pi} \mathrm{Im} 
        \sum_{m,n}
        \left[ 
        \mathsf{U}^{\alpha} \mathsf{S}_{\kvec}^{\phantom \dagger} \mathsf{T}_{\kvec}^{\phantom \dagger}
        \left( \int_{0}^{\infty} dt \, 
        \begin{pmatrix}
            (\mathsf{I} + \mathsf{n}_{\kvec}) e^{- i \mathsf{E}_{\kvec} t} & 0 \\
            0 & \mathsf{n}_{-\kvec} e^{i \mathsf{E}_{-\kvec} t}
        \end{pmatrix}
        \sin (\omega t)
        \right)
        \mathsf{T}^{\dagger}_{\kvec} \bar{\mathsf{S}}_{\kvec}^{\phantom \dagger} \bar{\mathsf{U}}^{\beta}
        \right]_{mn} \\
        = & -  \frac{2}{\pi} \mathrm{Im} 
        \sum_{m,n}
        \left[ 
            \mathsf{U}^{\alpha} \mathsf{S}_{\kvec}^{\phantom \dagger} \mathsf{T}_{\kvec}^{\phantom \dagger}
            \begin{pmatrix}
                (\mathsf{I + n}_{  \kvec}) \mathcal{G}_{\kvec} (\omega ; \eta) & 0 \\
                0 &  \mathsf{n}_{- \kvec} \mathcal{G}_{-\kvec}^{*} (\omega ; \eta)
            \end{pmatrix}
            \mathsf{T}^{\dagger}_{\kvec} \bar{\mathsf{S}}_{\kvec}^{\phantom \dagger}
            \bar{\mathsf{U}}^{\beta}
        \right]_{m,n}
    \end{split}
\end{equation}
where the function $\mathcal{G}_{\kvec}(\omega ; \eta)$ is defined as 
\begin{align*}
     \mathcal{G}_{\kvec}(\omega ; \eta)
     \equiv & \int_{0}^{\infty} dt \,  e^{- \ri E_{\kvec} t} \sin (\omega t) e^{-\eta |t|} 
     = \int_{0}^{\infty} dt \,  e^{-\ri E_{\kvec} t} \bigg( \frac{e^{\ri \omega t} - e^{- \ri \omega t}}{2\ri} \bigg) e^{-\eta t} \\
     = & \frac{1}{2 \ri} 
     \left( 
     \int_{0}^{\infty} dt \,
     e^{-(\eta + \ri(E_{\kvec} - \omega))t} - \int_{0}^{\infty} dt \, e^{-(\eta + \ri(E_{\kvec} + \omega))t} 
     \right) 
     = \frac{1}{2 \ri} \Big( \frac{1}{\eta + \ri(E_{\kvec} - \omega)} - \frac{1}{\eta + \ri(E_{\kvec} + \omega)} \Big) \\
     = & \frac{1}{2i} \big( g_{\eta}(\kvec, \omega) + g_{\eta} (\kvec, - \omega) \big).
\end{align*}
Thus, the spectral function in LSWT can be computed from the product of the matrices.
\begin{equation}
    \mathcal{A}^{\alpha\beta}(\kvec, \omega) 
    = - \frac{2}{\pi} \mathrm{Im} 
    \sum_{m,n}
    \left[ 
    \mathsf{U}^{\alpha} \mathsf{S}_{\kvec}^{\phantom \dagger} \mathsf{T}_{\kvec}^{\phantom \dagger}
    \mathcal{G} [ \mathsf{N}_{\kvec} ] (\omega ; \eta)
    \mathsf{T}^{\dagger}_{\kvec} \bar{\mathsf{S}}_{\kvec}^{\phantom \dagger} \bar{\mathsf{U}}^{\beta}
    \right]_{m,n},
\end{equation}
where the matrix $\mathcal{G} [ \mathsf{N}_{\kvec} ] (\omega ; \eta)$ is defined as 
\begin{equation}
    \mathcal{G} [ \mathsf{N}_{\kvec} ] (\omega ; \eta)
    = 
    \begin{pmatrix}
        (\mathsf{I} + \mathsf{n}_{\kvec}) \mathcal{G}_{\kvec} (\omega; \eta) & 0 \\
        0 & \mathsf{n}_{-\kvec} \mathcal{G}_{-\kvec}^{*} (\omega; \eta).
    \end{pmatrix}.
\end{equation}





\newpage
\section{Skyrmions, Topological Magnons, and Hall Effects}

\subsection{Skyrmion number}\label{subsec: Skyrmion number}

The skyrmion number itself does not require the LSWT, but introducing it here would be useful for a later purpose. 
ho in the lattice~\cite{BERG1981412} 
However,
The skyrmion number $Q_{\rm sk}$ quantifies the topological charge of a magnetic skyrmion in a lattice, computed as the sum of solid angles over triangular plaquettes ($\triangle$):
\begin{equation}
    Q_{\rm skyr} 
    = \frac{1}{4\pi} \sum_{\triangle} \chi_{\triangle}.
\end{equation}
For example, when the subsystem has four magnetic sublattices ($\mu = a, b, c, d$), the summation runs over the triangles $abc, acd, abd, bcd$.
Here, $\chi_{\triangle}$ is the solid angle subtended by normalized magnetization vectors $\mathbf{m}_i$, $\mathbf{m}_j$, and $\mathbf{m}_k$ at the vertices of a plaquette, where $\mathbf{m}_i = \mathbf{S}_i / S_i$ with $|\mathbf{m}_i| = 1$. The solid angle is given by:
\begin{equation}
    \chi_{\triangle} 
    = 2 \arctan 
    \left( 
    \frac{|\mathbf{m}_i \cdot (\mathbf{m}_j \times \mathbf{m}_k)|}{1 + \mathbf{m}_i \cdot \mathbf{m}_j + \mathbf{m}_j \cdot \mathbf{m}_k + \mathbf{m}_k \cdot \mathbf{m}_i} 
    \right).
\end{equation}





\subsection{Chern Number}\label{subsec: Chern number}

The Chern number of the $n$th band, $C_n$, is defined as the integral of the Berry curvature $\Omega_{n, \kvec}$ over the first Brillouin zone (FBZ):
\begin{equation}
    C_n = \frac{1}{2\pi} \int_{\text{FBZ}} \Omega_{n, \kvec} \, d^2\mathbf{k} 
    \approx \frac{1}{2\pi} \sum_{\kvec \in \text{FBZ}} \Omega_{n, \kvec} \Delta k_x \Delta k_y.
\end{equation}
The Berry curvature for the $n$th band is calculated as
\begin{equation}
    \Omega_{n, \kvec}^{\mu\nu} = 
    - 2 \, \mathrm{Im} \sum_{\substack{m = 1 \\ m \neq n}}^{2N}
    \frac{
        \big[ \mathsf{J} \mathsf{T}^{\dagger}_{\kvec} (\partial_{\mu} \mathsf{H}_{\kvec}) \mathsf{T}_{\kvec} \big]_{nm}
        \big[ \mathsf{J} \mathsf{T}^{\dagger}_{\kvec} (\partial_{\nu} \mathsf{H}_{\kvec}) \mathsf{T}_{\kvec} \big]_{mn}
    }{\big( (\mathsf{J} \mathsf{E}_{\kvec})_{nn} - (\mathsf{J} \mathsf{E}_{\kvec})_{mm} \big)^2},
\end{equation}
where $\mathsf{J}$, $\mathsf{H}_{\kvec}$, $\mathsf{T}_{\kvec}$, and $\mathsf{E}_{\kvec}$ are $2N \times 2N$ matrices, defined earlier
\begin{equation}
    \mathsf{J} = 
    \begin{pmatrix}
        \mathsf{I} & 0 \\ 0 & -\mathsf{I}
    \end{pmatrix}, 
    \qquad \text{and} \qquad
    \mathsf{T}^{\dagger}_{\kvec} \mathsf{H}_{\kvec} \mathsf{T}_{\kvec} = \mathsf{E}_{\kvec},
\end{equation}
with $\mathsf{T}_{\kvec}$ diagonalizing the Hamiltonian matrix $\mathsf{H}_{\kvec}$ to yield the energy matrix $\mathsf{E}_{\kvec}$.




\subsection{Thermal Hall Conductance}\label{subsec: Thermal Hall Conductance}

The thermal Hall conductivity (THC), while it's hard to observe in experiments, is related to the topology of magnon bands~\cite{PhysRevLett.128.117201}. It is computed using the Berry curvature as
\begin{equation}
    \kappa_{xy} = - \frac{k_B^2 T}{V} \sum_{\mathbf{k} \in \mathrm{FBZ}} \sum_{n=1}^{N} \left\{ c_2[g(\varepsilon_{n, \kvec})] - \frac{\pi^2}{3} \right\} \Omega_{n, \kvec},
\end{equation}
where $V$ is the system volume, $g(\varepsilon_{n, \kvec}) = [e^{\varepsilon_{n, \kvec}/k_B T} - 1]^{-1}$ is the Bose-Einstein distribution, and $c_2(x)$ is the Spence function, defined as
\begin{equation}
    c_2(x) 
    = \int_{0}^{x} dt \, \left( \ln \frac{1 + t}{t} \right)^2 
    = (1 + x) \left( \ln \frac{1 + x}{x} \right)^2 - (\ln x)^2 - 2 \mathrm{Li}_2(-x), 
    \quad \text{with} \quad 
    \mathrm{Li}_2(z) = - \int_{0}^{z} \frac{\ln (1 - t)}{t} \, dt.
\end{equation}



\newpage


\section{Example: Solving spin system with linear spin wave theory}

In this example, we solve a quadratic boson Hamiltonian. Consider the following Hamiltonian in momentum space, rewritten in Nambu form:
\begin{equation}
    H = \sum_{\mathbf{k}} A_{\mathbf{k}} a^{\dagger}_{\mathbf{k}} a_{\mathbf{k}} + \frac{1}{2} \Big( B_{\mathbf{k}} a^{\dagger}_{\mathbf{k}} a^{\dagger}_{-\mathbf{k}} + B_{\mathbf{k}}^{*} a_{\mathbf{k}} a_{-\mathbf{k}} \Big) 
    = \frac{1}{2} \sum_{\mathbf{k}}
    \begin{pmatrix}
        a^{\dagger}_{\mathbf{k}} & a_{-\mathbf{k}}
    \end{pmatrix}
    \begin{pmatrix}
        A_{\mathbf{k}} & B_{\mathbf{k}} \\
        B_{\mathbf{k}}^{*} & A_{\mathbf{k}}
    \end{pmatrix}
    \begin{pmatrix}
        a_{\mathbf{k}} \\
        a^{\dagger}_{-\mathbf{k}}
    \end{pmatrix},
\end{equation}
where $ B_{\mathbf{k}} $ is real without loss of generality. This assumption holds because we can redefine the bosonic operator with a phase factor $ a_{\mathbf{k}} \to a_{\mathbf{k}} e^{-i\phi} $, where $ B_{\mathbf{k}} = e^{i\phi} |B_{\mathbf{k}}| $, making $ B_{\mathbf{k}} $ real by choosing an appropriate $ \phi $.
Note that the condition to be Hamiltonian density positive is 
\begin{equation*}
    \lambda(\mathsf{H}_{\kvec}) = A_{\kvec} \pm B_{\kvec} \ge 0, 
    \qquad \rightarrow \qquad
    | A_{\kvec} | \ge | B_{\kvec} |.
\end{equation*}
This condition corresponds to the positivity of the the magnon energy $\omega_{\kvec} = \sqrt{A_{\kvec}^{2} - B_{\kvec}^{2}}$.

To diagonalize the Hamiltonian, we introduce the Bogoliubov transformation:
\begin{equation}
    \begin{pmatrix}
        a_{\mathbf{k}} \\
        a^{\dagger}_{-\mathbf{k}}
    \end{pmatrix}
    =
    \begin{pmatrix}
        \cosh \theta_{\mathbf{k}} & \sinh \theta_{-\mathbf{k}} \\
        \sinh \theta_{\mathbf{k}} & \cosh \theta_{-\mathbf{k}}
    \end{pmatrix}
    \begin{pmatrix}
        \alpha_{\mathbf{k}} \\
        \alpha^{\dagger}_{-\mathbf{k}}
    \end{pmatrix}
\end{equation}
where $ \theta_{\mathbf{k}} $ is real and even under $ \mathbf{k} \to -\mathbf{k} $, so $ \theta_{\mathbf{k}} = \theta_{-\mathbf{k}} $. The bosonic commutation relations $ [a_{\mathbf{k}}, a^{\dagger}_{\mathbf{k}}] = 1 $ and $ [\alpha_{\mathbf{k}}, \alpha^{\dagger}_{\mathbf{k}}] = 1 $ must be preserved. Define the transformation matrix as:
\begin{equation}
    \mathsf{T}_{\mathbf{k}} = 
    \begin{pmatrix}
        \cosh \theta_{\mathbf{k}} & \sinh \theta_{\mathbf{k}} \\
        \sinh \theta_{\mathbf{k}} & \cosh \theta_{\mathbf{k}}
    \end{pmatrix}
\end{equation}
Check the commutation relation for the new operators:
\begin{align}
    [\alpha_{\mathbf{k}}, \alpha^{\dagger}_{\mathbf{k}}] 
    &= [\cosh \theta_{\mathbf{k}} a_{\mathbf{k}} + \sinh \theta_{\mathbf{k}} a^{\dagger}_{-\mathbf{k}}, \cosh \theta_{\mathbf{k}} a^{\dagger}_{\mathbf{k}} + \sinh \theta_{\mathbf{k}} a_{-\mathbf{k}}]   \\
    &= \cosh^2 \theta_{\mathbf{k}} [a_{\mathbf{k}}, a^{\dagger}_{\mathbf{k}}] + \sinh^2 \theta_{\mathbf{k}} [a^{\dagger}_{-\mathbf{k}}, a_{-\mathbf{k}}]   \\
    &= \cosh^2 \theta_{\mathbf{k}} \cdot 1 + \sinh^2 \theta_{\mathbf{k}} \cdot (-1)   \\
    &= \cosh^2 \theta_{\mathbf{k}} - \sinh^2 \theta_{\mathbf{k}} = 1
\end{align}
This holds due to the hyperbolic identity, confirming that the transformation is canonical.

\subsection*{Diagonalizing the Hamiltonian}

Substitute the transformation into the Hamiltonian. First, the conjugate transformation is:
\begin{equation}
    \begin{pmatrix}
        a^{\dagger}_{\mathbf{k}} & a_{-\mathbf{k}}
    \end{pmatrix}
    =
    \begin{pmatrix}
        \alpha^{\dagger}_{\mathbf{k}} & \alpha_{-\mathbf{k}}
    \end{pmatrix}
    \begin{pmatrix}
        \cosh \theta_{\mathbf{k}} & \sinh \theta_{\mathbf{k}} \\
        \sinh \theta_{\mathbf{k}} & \cosh \theta_{\mathbf{k}}
    \end{pmatrix}
\end{equation}
Thus, the Hamiltonian becomes:
\begin{equation}
    H = \frac{1}{2} \sum_{\mathbf{k}}
    \begin{pmatrix}
        \alpha^{\dagger}_{\mathbf{k}} & \alpha_{-\mathbf{k}}
    \end{pmatrix}
    T_{\mathbf{k}}
    \begin{pmatrix}
        A_{\mathbf{k}} & B_{\mathbf{k}} \\
        B_{\mathbf{k}}^{*} & A_{\mathbf{k}}
    \end{pmatrix}
    T_{\mathbf{k}}
    \begin{pmatrix}
        \alpha_{\mathbf{k}} \\
        \alpha^{\dagger}_{-\mathbf{k}}
    \end{pmatrix}
\end{equation}
Define the matrix to be diagonalized as:
\begin{equation}
    \mathsf{E}_{\mathbf{k}} = \mathsf{T}_{\mathbf{k}}^{\dagger}
    \begin{pmatrix}
        A_{\mathbf{k}} & B_{\mathbf{k}} \\
        B_{\mathbf{k}}^{*} & A_{\mathbf{k}}
    \end{pmatrix}
    \mathsf{T}_{\mathbf{k}}
    =
    \begin{pmatrix}
        \cosh \theta_{\mathbf{k}} & \sinh \theta_{\mathbf{k}} \\
        \sinh \theta_{\mathbf{k}} & \cosh \theta_{\mathbf{k}}
    \end{pmatrix}
    \begin{pmatrix}
        A_{\mathbf{k}} & B_{\mathbf{k}} \\
        B_{\mathbf{k}}^{*} & A_{\mathbf{k}}
    \end{pmatrix}
    \begin{pmatrix}
        \cosh \theta_{\mathbf{k}} & \sinh \theta_{\mathbf{k}} \\
        \sinh \theta_{\mathbf{k}} & \cosh \theta_{\mathbf{k}}
    \end{pmatrix}
\end{equation}
Since $ B_{\mathbf{k}} $ is real, $ B_{\mathbf{k}}^{*} = B_{\mathbf{k}} $.
Compute $ M_{\mathbf{k}} $:
\begin{equation}
    \mathsf{E}_{\mathbf{k}}
    = 
    \begin{pmatrix}
        A_{\mathbf{k}} \cosh 2\theta_{\mathbf{k}} + B_{\mathbf{k}} \sinh 2\theta_{\mathbf{k}} 
        & A_{\mathbf{k}} \sinh 2\theta_{\mathbf{k}} + B_{\mathbf{k}} \cosh 2\theta_{\mathbf{k}}\\
        A_{\mathbf{k}} \sinh 2\theta_{\mathbf{k}} + B_{\mathbf{k}} \cosh 2\theta_{\mathbf{k}} 
        & A_{\mathbf{k}} \cosh 2\theta_{\mathbf{k}} + B_{\mathbf{k}} \sinh 2\theta_{\mathbf{k}} 
    \end{pmatrix}
    = 
    \begin{pmatrix}
        \omega_{\kvec}
        & 0 \\ 0 & 
        \omega_{\kvec}
    \end{pmatrix}
\end{equation}
Thus, we have the diagonalization condition
\begin{equation}
    [\mathsf{E}_{\mathbf{k}}]_{12} = A_{\mathbf{k}} \sinh 2\theta_{\mathbf{k}} + B_{\mathbf{k}} \cosh 2\theta_{\mathbf{k}} = 0,
    \quad \rightarrow \quad
    \tanh 2\theta_{\mathbf{k}} = -\frac{B_{\mathbf{k}}}{A_{\mathbf{k}}}
\end{equation}
The diagonal element $M_{\kvec}^{(11)}$ gives the eigenenergy:
\begin{equation}
    \omega_{\mathbf{k}} 
    = A_{\mathbf{k}} \cosh 2\theta_{\mathbf{k}} + B_{\mathbf{k}} \sinh 2\theta_{\mathbf{k}} 
    = A_{\mathbf{k}} \cdot \frac{A_{\mathbf{k}}}{\sqrt{A_{\mathbf{k}}^2 - B_{\mathbf{k}}^2}} + B_{\mathbf{k}} \cdot \left(-\frac{B_{\mathbf{k}}}{\sqrt{A_{\mathbf{k}}^2 - B_{\mathbf{k}}^2}}\right) 
    = \sqrt{A_{\mathbf{k}}^2 - B_{\mathbf{k}}^2}.
\end{equation}
To express this in terms of $ A_{\mathbf{k}} $ and $ B_{\mathbf{k}} $, use $ \cosh 2\theta_{\mathbf{k}} = \frac{1}{\sqrt{1 - \tanh^2 2\theta_{\mathbf{k}}}} $ and $ \sinh 2\theta_{\mathbf{k}} = \tanh 2\theta_{\mathbf{k}} \cosh 2\theta_{\mathbf{k}} $:
\begin{equation*}
    \cosh 2\theta_{\mathbf{k}} = \frac{1}{\sqrt{1 - \left(-B_{\mathbf{k}} / A_{\mathbf{k}} \right)^2}} = \frac{A_{\mathbf{k}}}{\omega_{\kvec}},  
    \qquad 
    \sinh 2\theta_{\mathbf{k}} = -\frac{B_{\mathbf{k}}}{A_{\mathbf{k}}} \cdot \frac{A_{\mathbf{k}}}{\sqrt{A_{\mathbf{k}}^2 - B_{\mathbf{k}}^2}} = -\frac{B_{\mathbf{k}}}{\omega_{\kvec}}.
\end{equation*}


\noindent\textbf{Number Expectation in Ground State:}
\begin{align}
    \langle n_{\kvec} \rangle 
    = & \langle a_{\kvec}^{\dagger} a_{\kvec} \rangle 
    =  \langle 
    ( \cosh \theta_{\kvec} \alpha_{\kvec}^{\dagger} + \sinh \theta_{\kvec} \alpha_{-\kvec} ) 
    ( \cosh \theta_{\kvec} \alpha_{\kvec} + \sinh \theta_{\kvec} \alpha_{-\kvec}^{\dagger} ) 
    \rangle \\
    = & \cosh^{2} \theta_{\kvec} \langle  \alpha_{\kvec}^{\dagger} \alpha_{\kvec} \rangle 
    + \sinh^{2} \theta_{\kvec} \langle \alpha_{-\kvec} \alpha_{-\kvec}^{\dagger} \rangle \
    + \sinh \theta_{\kvec} \cosh \theta_{\kvec} 
    \big(\langle \alpha_{-\kvec}^{\dagger} \alpha_{\kvec}^{\dagger} \rangle + \langle \alpha_{\kvec} \alpha_{-\kvec} \rangle \big) \\
    = & \sinh^{2} \theta_{\kvec} 
    = \frac{\cosh 2 \theta_{\kvec} - 1}{2} 
    = \frac{A_{\kvec} / \omega_{\kvec} - 1}{2}.
\end{align}

\begin{itemize}
    \item If $\omega_{\kvec} = \sqrt{A^{2}_{\kvec} - |B_{\kvec}|^{2}} \rightarrow 0^{+}$, then the number 
    \begin{equation}
        \langle a_{k}^{\dagger} a_{k} \rangle = |v_k|^2 = \frac{A_{\kvec}/\omega_{\kvec} - 1}{2} \approx \frac{A_{\kvec}}{\omega_{\kvec}} \rightarrow \infty
    \end{equation}
    \item  If $B_{\kvec} = 0$:
    \begin{equation}
        \langle n_{\kvec} \rangle = 0
    \end{equation}
    \item If $A_{\kvec} \rightarrow A_{\kvec} + \mu$, which corresponds to adding on-site chemical potential amount to $\mu \sum_{\kvec} a_{\kvec}^{\dagger}a_{\kvec} = \frac{\mu}{2} \sum_{\kvec} (a_{\kvec}^{\dagger}a_{\kvec} + a_{-\kvec}a^{\dagger}_{-\kvec} )$ is the quantum correction
    \begin{align*}
        E_{\rm qu}^{(\mu)}
        = & \frac{1}{2} \sum_{\kvec} \Big( \sum_{\mu} E_{\kvec, \mu} - \Tr[\mathsf{A}_{\kvec}]\Big)
        = \frac{1}{2}\sum_{\kvec}  \Big(  ((A_{\kvec} + \mu)^{2} - B_{\kvec}^{2})^{1/2} - (A_{\kvec} + \mu) \Big) \\
        = & \frac{1}{2} \sum_{\kvec} \big( (A_{\kvec}^{2} - B_{\kvec})^{1/2} - A_{\kvec} \big) 
        + \Big( \frac{A_{\kvec}}{\omega_{\kvec}} - 1\Big) \mu 
        + \frac{B_{\kvec}^{2}}{2\omega^{3}_{\kvec}} \mu^{2} + \cdots  \\
        = & \sum_{\kvec} \frac{\omega_{\kvec}}{2} + \langle a_{\kvec}^{\dagger} a_{\kvec} \rangle \mu +  \frac{\langle a_{\kvec}^{\dagger} a_{\kvec}^{\dagger} a_{\kvec} a_{\kvec} \rangle}{\omega_{\kvec}} \mu^{2} + \cdots.
    \end{align*}
    \begin{align*}
        \langle n_{\kvec}^{(\mu)} \rangle - \langle n_{\kvec}^{(0)} \rangle 
        = & \frac{A_{\kvec} + \mu - \sqrt{(A_{\kvec} + \mu)^{2} - B_{\kvec}^{2}}}{2 \sqrt{(A_{\kvec} + \mu)^{2} - B_{\kvec}^{2}} } 
        - \frac{A_{\kvec} - \sqrt{(A_{\kvec})^{2} - B_{\kvec}^{2}}}{2 \sqrt{(A_{\kvec})^{2} - B_{\kvec}^{2}} } \\
        = & - \frac{B_{\kvec}^{2}}{2 \omega^{3}_{\kvec}} \mu - \frac{3}{4} \frac{A_{\kvec}B_{\kvec}^{2}}{\omega_{\kvec}^{5}} \mu^{2} + \cdots
    \end{align*}
\end{itemize}


\bibliography{Setting/citation}


\end{document}



